\section{Adelic Riemann--Roch: the repartition cokernel (RR.A)}
\label{mk-chap:RiemannRoch_Adelic}

\section{The two-chart model}
\label{mk-sec:RiemannRoch_Adelic_intro}
The concrete repartition presentation is inspired by Papaioannou \cite{papaioannou-rr}.

Let \(C\) be a smooth proper geometrically integral curve over a field
\(k\), and let \(K=k(C)\).  The divisor theory of
\ref{mk-chap:RiemannRoch_WeilDivisor} assigns to every prime divisor \(P\)
an order \(v_P\colon K^\times\to\mathbb Z\).  The Mayer--Vietoris theory
of \ref{mk-chap:Cohomology_MayerVietoris} computes the first cohomology of
the structure sheaf from two affine charts.

\paragraph{The concrete-repartition idea.}
Classically (Weil), for the function field \(K = k(C)\) one forms the
\emph{adele} (repartition) space
\(\mathbb{A}_K = \{(x_P)_P \in \prod_P K : v_P(x_P) \ge 0 \text{ for
almost all } P\}\), its bounded subspaces
\(\mathbb{A}_K(D) = \{(x_P) : v_P(x_P) \ge -D(P)\ \forall P\}\), and the
two Riemann--Roch spaces
\[
  L(D) = K \cap \mathbb{A}_K(D)
        = \{ f \in K : \operatorname{div} f + D \ge 0\},
  \qquad
  H^1(D) = \mathbb{A}_K / (\mathbb{A}_K(D) + K).
\]
The genus is \(g=\dim_k H^1(0)\), and the index of speciality is
\(i(D)=\dim_k H^1(D)\).  A two-affine cover gives a finite presentation
of the same cohomology without first constructing the restricted product.

\paragraph{The \(2\)-affine-cover cokernel.}
Fix affine opens \(U_0,U_1\) covering \(C\), with affine intersection, and
write
\[
  A_i = \Gamma(U_i, \struct{C}), \qquad
  A_{01} = \Gamma(U_0 \cap U_1, \struct{C}).
\]
The degree-\(1\) Čech cohomology of the two-element cover is the cokernel
of the difference map \((f_0, f_1) \mapsto f_0|_{01} - f_1|_{01}\):
\[
  \check{H}^1 \;=\; A_{01} \big/ (A_0 + A_1).
\]
For a Weil divisor \(D\) the \emph{twisted} version replaces each ring of
sections by the Riemann--Roch sheaf of sections
\(\Gamma(U, \struct{C}(D)) = \{ f \in K : v_P(f) \ge -D(P) \text{ for all
prime divisors } P \in U\}\), giving
\[
  L(D) = \Gamma(U_0, \struct{C}(D)) \cap \Gamma(U_1, \struct{C}(D)).
\]
\[
  \check{H}^1(D) = \Gamma(U_0 \cap U_1, \struct{C}(D)) \big/
    \bigl(\Gamma(U_0, \struct{C}(D)) + \Gamma(U_1, \struct{C}(D))\bigr).
\]
The bridge of \ref{mk-sec:RiemannRoch_Adelic_cokernel} identifies
\(\check H^1(0)\) with \(H^1(C,\struct C)\).

\paragraph{Residue-degree bookkeeping.}
Because \(k\) need not be algebraically closed, every place carries a
\emph{residue degree}. For a prime divisor \(P\) with residue field
\(\kappa(P) = \struct{C, P}/\mathfrak{m}_P\), set
\[
  \deg P \;=\; [\kappa(P) : k] \;=\; \dim_k \kappa(P) \;<\; \infty,
  \qquad
  \deg_k D \;=\; \sum_P D(P)\,\deg P .
\]
Over an algebraically closed field this agrees with the coefficient degree
of \ref{mk-def:divisor_degree}; over a general field it does not.  In this
chapter \(\deg D\) always means the residue-weighted degree \(\deg_kD\).

We write \(\ell(D) = \dim_k L(D)\),
\(i(D) = \dim_k \check{H}^1(D)\), and
\(\chi(D) = \ell(D) - i(D)\).

\section{Valuations and divisor sections}
\label{mk-sec:RiemannRoch_Adelic_substrate}

\begin{proposition}
[Affine chart is a Dedekind domain with fraction field \(K\)]
  \label{mk-thm:adelic_chartRing_isDedekindDomain}
  \lean{AlgebraicGeometry.Adelic.chartRing_isFractionRing,
    AlgebraicGeometry.Adelic.isDedekindDomain_chart}
  \leanok
  \dcref{custom}

  \uses{mk-def:isRegularInCodimensionOne}
  Let \(U \subseteq C\) be a nonempty affine open with coordinate ring
  \(A_U = \Gamma(U, \struct{C})\).  The canonical map
  \(A_U\hookrightarrow K\) exhibits \(K\) as its fraction field.  Under
  the hypothesis that every nonempty affine chart of \(C\) is Dedekind,
  \(A_U\) is a Dedekind domain.
\end{proposition}
\begin{proof}
  \leanok
  The fraction-field assertion is localization at the generic point.
  The Dedekind assertion is the specialization of the stated chart
  hypothesis to \(U\).
\end{proof}

\begin{definition}[Point valuation]
  \label{mk-def:adelic_pointValuation}
  \lean{AlgebraicGeometry.Adelic.pointValuation}
  \leanok
  \source{papaioannou-algebraic-rr:page-0003,
    papaioannou-algebraic-rr:page-0004}

  \uses{mk-def:prime_divisor, mk-def:isRegularInCodimensionOne}
  \dcref{custom}

  For a prime divisor \(P\), the discrete valuation ring
  \(\struct{C,P}\) defines its adic valuation
  \[
    \nu_P\colon K\longrightarrow\mathbb Z_{\mathrm m}^{0}.
  \]
\end{definition}

\begin{lemma}[Order is the additive normalization of the point valuation]
  \label{mk-lem:adelic_order_eq_neg_log_pointValuation}
  \lean{AlgebraicGeometry.Adelic.order_eq_neg_log_pointValuation}
  \leanok
  \dcref{custom}

  \uses{mk-def:adelic_pointValuation, mk-def:order_at_point}

  For every \(f\in K\),
  \[
    \operatorname{ord}_P(f)=-\log\bigl(\nu_P(f)\bigr).
  \]
\end{lemma}
\begin{proof}
  \leanok
  This is the comparison between the height-one-stalk adic valuation and the
  \(\operatorname{ordFrac}\) normalization used in
  \ref{mk-def:order_at_point}.
\end{proof}

\begin{definition}[Residue field and residue degree]
  \label{mk-def:adelic_residue_degree}
  \uses{mk-thm:adelic_chartRing_isDedekindDomain, mk-def:adelic_pointValuation}
  \dcref{custom}

  If \(P\) lies in a nonempty affine open \(U\), it determines a height-one
  prime \(\mathfrak p_P\) of the Dedekind chart
  \(A_U=\Gamma(U,\struct C)\).  Its residue field and residue degree are
  \[
    \kappa(P)=A_U/\mathfrak p_P
      =\struct{C,P}/\mathfrak m_P,
    \qquad
    \deg P=[\kappa(P):k]=\dim_k\kappa(P).
  \]
  The stalk \(\struct{C, P}\) is the localisation \((A_U)_{\mathfrak p_P}\)
  because \(P\) lies in \(U\).  The definition is independent of the chart,
  and \(\deg P\) is finite by Zariski's lemma because \(P\) is a closed point
  of the finite-type \(k\)-scheme \(C\).
\end{definition}

\begin{definition}[Residue-weighted degree]
  \label{mk-def:adelic_divisor_degree}
  \uses{mk-def:adelic_residue_degree, mk-def:codim1_cycles}
  \dcref{custom}

  For a Weil divisor \(D=\sum_P n_P[P]\) on a curve over \(k\), define its
  residue-weighted degree by
  \[
    \deg_k(D)=\sum_P n_P[\kappa(P):k].
  \]
  The sum is finite because a Weil divisor has finite support.  If \(k\) is
  algebraically closed, every closed point has residue field \(k\), so
  \(\deg_k\) agrees with the coefficient degree of
  \ref{mk-def:divisor_degree}.
\end{definition}

\begin{theorem}[Weighted degree of a principal divisor]
  \label{mk-thm:adelic_principal_degree_zero}
  \uses{mk-def:adelic_divisor_degree,
    mk-thm:adelic_exists_finiteMorphismToP1}
  \dcref{custom}
  For every \(f\in K^\times\),
  \[
    \deg_k\operatorname{div}(f)=0.
  \]
\end{theorem}
\begin{proof}
  A constant function has zero order everywhere.  Otherwise \(f\) defines
  a finite morphism \(\varphi_f\colon C\to\mathbb P^1_k\), and
  \(\operatorname{div}(f)=\varphi_f^*([0]-[\infty])\).  The
  residue-weighted degree of a finite pullback is multiplied by
  \(\deg\varphi_f\); since both \([0]\) and \([\infty]\) have degree one,
  the displayed divisor has degree zero.
\end{proof}

\begin{definition}
[Sections bounded by a divisor; the Riemann--Roch space]
  \label{mk-def:adelic_sectionOfDivisor}
  \lean{AlgebraicGeometry.Adelic.sectionOfDivisor,
    AlgebraicGeometry.Adelic.linearSystem,
    AlgebraicGeometry.Adelic.linearSystem_eq_inf}
  \leanok
  \source{papaioannou-algebraic-rr:page-0007}

  \uses{mk-def:adelic_pointValuation, mk-def:codim1_cycles, mk-lem:order_mul_of_ne_zero}
  \dcref{custom}
  For an open \(U \subseteq C\) and a Weil divisor \(D\)
  (\ref{mk-def:codim1_cycles}) define the additive subgroup of \(K\)
  \[
    \Gamma(U, \struct{C}(D)) \;=\;
      \{ f \in K : f = 0 \text{ or } v_P(f) \ge -D(P) \text{ for every
         prime divisor } P \text{ with point in } U\}
  \]
  (the disjunct \(f = 0\) encodes \(v_P(0) = +\infty\)).  This is an
  additive subgroup of \(K\).  Its global instance gives the
  \emph{Riemann--Roch space}
  \(L(D) = \Gamma(U_0, \struct{C}(D)) \cap \Gamma(U_1, \struct{C}(D))\),
  equal to \(\{ f \in K : \operatorname{div} f + D \ge 0\}\) since
  \(U_0 \cup U_1 = C\).
\end{definition}
\begin{proof}
  \leanok
  The set contains \(0\)
  (\(v_P(0) = +\infty \ge -D(P)\)). It is closed under addition by the
  ultrametric inequality of the discrete valuation
  \(v_P(f + h) \ge \min(v_P f, v_P h) \ge -D(P)\). It is closed under
  negation because \(v_P(-f)=v_P(f)\).
  The intersection identity for \(L(D)\) holds because a prime divisor of
  \(C\) has its point in \(U_0\) or in \(U_1\); imposing \(v_P(f) \ge
  -D(P)\) on the union of the two chart conditions is imposing it at
  every prime divisor, i.e.\ \(\operatorname{div} f + D \ge 0\).
\end{proof}

\begin{lemma}
[Finiteness of the non-integral places on a Dedekind chart]
  \label{mk-thm:adelic_finite_support_order}
  \lean{AlgebraicGeometry.Adelic.finite_places_ne_one}
  \leanok
  \dcref{custom}

  \uses{mk-thm:adelic_chartRing_isDedekindDomain}
  Let \(R\) be a Dedekind domain with fraction field \(K\). For a nonzero
  \(k \in K\) the set \(\{ v \in \operatorname{Spec}^1 R : v(k) \ne 1\}\) of
  height-one places at which \(k\) is a zero or a pole is finite.

  This is the affine-chart core of the finite-support statement: a curve is
  covered by finitely many Dedekind charts, and applying the lemma on each
  chart shows that for \(f \in K^\times\) only finitely many prime divisors
  have \(v_P(f) \ne 0\), so \(\operatorname{div} f = \sum_P v_P(f)\,P\) is a
  well-defined Weil divisor. The general scheme-level assembly (finite affine
  cover plus the height-one/minimal-primes bound) is
  \ref{mk-lem:rationalMap_order_finite_support}.
\end{lemma}
\begin{proof}
  \leanok
  Apply the finiteness of the support of a nonzero fractional ideal in a
  Dedekind domain to both \(k\) and \(k^{-1}\). A place \(v\) with
  \(v(k) \ne 1\) satisfies either \(v(k) > 1\) (a zero of \(k\)) or
  \(v(k) < 1\), equivalently \(v(k^{-1}) > 1\) (a pole of \(k\)); each of
  the two sets of such places is finite, and their union contains every
  place with \(v(k) \ne 1\).
\end{proof}

\section{The cover cokernel and first cohomology}
\label{mk-sec:RiemannRoch_Adelic_cokernel}

\begin{definition}
[The two-cover cokernel of a sheaf]
  \label{mk-def:adelic_coverH1}
  \lean{AlgebraicGeometry.Scheme.AffineCoverMVSquare.H1Cok}
  \leanok
  \dcref{custom}

  \uses{mk-def:Scheme_AffineCoverMVSquare}
  For a sheaf \(F\) of \(k\)-modules and an affine Mayer--Vietoris square,
  define
  \[
    \operatorname{H1Cok}_S(F)
      =\Gamma(U_0\cap U_1,F)/
        \operatorname{im}\bigl(\Gamma(U_0,F)\oplus\Gamma(U_1,F)
          \xrightarrow{\,f_0-f_1\,}\Gamma(U_0\cap U_1,F)\bigr).
  \]
\end{definition}

\begin{definition}[The function-field twist cokernel]
  \label{mk-def:adelic_function_field_H1}
  \lean{AlgebraicGeometry.Adelic.H1}
  \leanok
  \dcref{custom}

  \uses{mk-def:adelic_sectionOfDivisor}
  Independently of a sheaf structure, the divisor section subgroups define
  the additive quotient
  \[
    \check H^1_{\mathrm{ff}}(D)
      =\Gamma(U_0\cap U_1,\struct C(D))/
        \bigl(\Gamma(U_0,\struct C(D))+\Gamma(U_1,\struct C(D))\bigr).
  \]
\end{definition}

\begin{proposition}[Comparison at the zero divisor]
  \label{mk-thm:adelic_function_field_H1_zero_compare}
  \dcref{custom}
  \uses{mk-def:adelic_coverH1, mk-def:adelic_function_field_H1}
  There is a canonical isomorphism
  \[
    \check H^1_{\mathrm{ff}}(0)
      \simeq_k\operatorname{H1Cok}_S(\struct C).
  \]
\end{proposition}
\begin{proof}
  On every nonempty affine open of the integral normal curve, a rational
  function is regular exactly when it has nonnegative order at every
  closed point of that open.  This identifies each term and each
  restriction map in the two difference complexes, hence their cokernels.
\end{proof}

\begin{lemma}
[The unnormalized coface differential is the alternating coface sum]
  \label{mk-lem:adelic_alternatingCofaceMapComplex_objD}
  \lean{AlgebraicGeometry.Scheme.alternatingCofaceMapComplex_objD}
  \leanok
  \dcref{custom}
  For a cosimplicial object \(Y\) in a preadditive category and \(n \ge
  0\), the degree-\(n\) differential of the associated (unnormalized)
  alternating coface cochain complex is the alternating sum of coface
  maps,
  \[
    d^n \;=\; \sum_{i = 0}^{n+1} (-1)^i\, Y(\delta_i) \colon Y^n \to
    Y^{n+1}.
  \]
\end{lemma}
\begin{proof}
  \leanok
  This is the cosimplicial dual of the standard alternating-face-map
  formula for the (co)chain complex of a simplicial object: the
  differential of the alternating coface complex is defined as this
  alternating sum, so the identity is definitional unfolding.
\end{proof}

\begin{definition}
[The Čech cosimplicial object of a cover]
  \label{mk-def:adelic_cechCosimplicial}
  \lean{AlgebraicGeometry.Scheme.cechCosimplicial,
    AlgebraicGeometry.Scheme.cechCochain_eq}
  \leanok
  \dcref{custom}
  \uses{mk-def:adelic_coverH1, mk-def:Scheme_cechCohomology}
  For a family of opens \(\mathcal{U} = (U_i)_{i \in \iota}\) of \(C\)
  and a sheaf \(F\) of \(k\)-modules, the \emph{Čech cosimplicial object}
  \(Y_{\mathcal{U}, F}\) is given in degree \(n\) by
  \[
    Y_{\mathcal U,F}^n=
      \prod_{j \colon \mathrm{Fin}(n+1) \to \iota}
        \Gamma\Bigl(\bigsqcap_a U_{j(a)},F\Bigr).
  \]
  Its coface maps \(\delta_{i_0} \colon Y^n \to Y^{n+1}\) are induced by
  restriction along the shrinking of the
  intersection when an index is duplicated. The Čech cochain complex is
  by definition the alternating coface complex of \(Y_{\mathcal{U},
  F}\):
  \(\mathrm{cechCochain}\, C\, F\, \mathcal{U} =\) the alternating coface
  complex of \(Y_{\mathcal{U}, F}\)
  (\ref{mk-lem:adelic_alternatingCofaceMapComplex_objD}).
\end{definition}
\begin{proof}
  \leanok
  Immediate from the definition of \(\mathrm{cechCochain}\) as the
  alternating coface complex associated to the cosimplicial object
  obtained by evaluating the Čech nerve of \(\mathcal{U}\) against \(F\).
\end{proof}

\begin{lemma}
[Cofaces are restriction maps]
  \label{mk-lem:adelic_cech_coface_restriction}
  \lean{AlgebraicGeometry.Scheme.cechCosimplicial_δ_π,
    AlgebraicGeometry.Scheme.prodOpens_eq_iInf,
    AlgebraicGeometry.Scheme.prodOpens_δ_le,
    AlgebraicGeometry.Scheme.cechCosimplicial_δ_π_restrict}
  \leanok
  \dcref{custom}
  \uses{mk-def:adelic_cechCosimplicial}
  For every finite index family \(j \colon \mathrm{Fin}\, m \to \iota\),
  the categorical product \(\prod_a U_{j(a)}\) equals the intersection
  \(\bigsqcap_a U_{j(a)}\) of opens of \(C\). Consequently, for
  \(n \ge 0\), \(i_0 \in \mathrm{Fin}(n+2)\), and
  \(j \colon \mathrm{Fin}(n+2) \to \iota\), the coface map \(\delta_{i_0}
  \colon Y_{\mathcal{U},F}^n \to Y_{\mathcal{U},F}^{n+1}\), projected to
  the \(j\)-indexed factor of the target, is the restriction map
  \[
    \Gamma\Bigl(\bigsqcap_a U_{j(\delta_{i_0}(a))}, F\Bigr)
      \longrightarrow
    \Gamma\Bigl(\bigsqcap_a U_{j(a)}, F\Bigr)
  \]
  along the inclusion
  \(\bigsqcap_a U_{j(a)} \subseteq \bigsqcap_a U_{j(\delta_{i_0}(a))}\)
  (dropping the index \(i_0\) enlarges the intersection).
\end{lemma}
\begin{proof}
  \leanok
  The abstract categorical product of a finite family of opens, computed
  in the poset \(\mathrm{Opens}\, C\), is its infimum, by the universal
  property of both constructions. Under this identification the coface
  map's action on the \(j\)-indexed factor --- reindexing the source
  factor along \(j \circ \delta_{i_0}\) followed by the canonical
  comparison map between the two products --- is exactly the restriction
  map along the inclusion of the larger-index intersection into the
  smaller-index one; since \(\mathrm{Opens}\, C\) is a poset, this
  comparison map is the unique restriction along that inclusion.
\end{proof}

\begin{lemma}
[The degree-\(0\) and degree-\(1\) Čech differentials]
  \label{mk-lem:adelic_cech_d01_d12_formula}
  \lean{AlgebraicGeometry.Scheme.cechCochain_d01_eq,
    AlgebraicGeometry.Scheme.cechCochain_d12_eq}
  \leanok
  \dcref{custom}
  \uses{mk-lem:adelic_alternatingCofaceMapComplex_objD,
    mk-lem:adelic_cech_coface_restriction}
  In terms of the coface maps of \(Y_{\mathcal{U}, F}\), the Čech
  differentials in degrees \(0\) and \(1\) are
  \[
    d^0 = \delta_0 - \delta_1, \qquad\qquad
    d^1 = \delta_0 - \delta_1 + \delta_2.
  \]
\end{lemma}
\begin{proof}
  \leanok
  Immediate from \ref{mk-lem:adelic_alternatingCofaceMapComplex_objD} at
  \(n = 0\) and \(n = 1\), expanding the alternating sum over
  \(\mathrm{Fin}\, 2\) and \(\mathrm{Fin}\, 3\) respectively.
\end{proof}

\begin{proposition}
[Two-element Čech \(H^1\) is the cokernel]
  \label{mk-thm:adelic_cechH1_eq_cokernel}
  \lean{AlgebraicGeometry.Scheme.AffineCoverMVSquare.cechCohomologyOneEquivH1Cok}
  \leanok
  \dcref{custom}
  \uses{mk-def:adelic_coverH1, mk-def:Scheme_cechCohomology, mk-def:Scheme_toModuleKSheaf}
  For any sheaf \(F\) of \(k\)-modules on \(C\) and any affine
  Mayer--Vietoris square \(S\) with cover family \(\mathcal U = \{U_0,
  U_1\}\), there is a \(k\)-linear isomorphism
  \[
    \mathrm{cechCohomology}\, C\, F\, \mathcal U\, 1
    \;\simeq_k\; \Gamma(U_0 \cap U_1, F) \big/
      \bigl(\Gamma(U_0, F) + \Gamma(U_1, F)\bigr),
  \]
  the image of \(\Gamma(U_0, F) \oplus \Gamma(U_1, F)\) under
  \((f_0, f_1) \mapsto f_0|_{01} - f_1|_{01}\) being identified with the
  sum \(\Gamma(U_0, F) + \Gamma(U_1, F)\) inside \(\Gamma(U_0 \cap U_1,
  F)\). Specialising \(F = \toModuleKSheaf C\) gives
  \[
    \mathrm{cechCohomology}\, C\, (\toModuleKSheaf C)\, S.\cover\, 1
      \;\simeq_k\; \check{H}^1 = A_{01}/(A_0 + A_1).
  \]
\end{proposition}
\begin{proof}
  \leanok
  \uses{mk-lem:adelic_cech_d01_d12_formula, mk-lem:adelic_cech_coface_restriction}
  Unfold \(\mathrm{cechCohomology}\, C\, F\, \mathcal U\, 1\) as the
  homology at the middle term of the short complex
  \(C^0 \xrightarrow{d^0} C^1 \xrightarrow{d^1} C^2\) of the Čech
  cochain complex, i.e.\ \(\ker d^1 / \operatorname{im}(C^0 \to \ker
  d^1)\). For the two-element index set \(\mathcal U = \{U_0, U_1\}\)
  every triple of indices \(x \colon \mathrm{Fin}\,3 \to \{0,1\}\)
  repeats a value, so by \ref{mk-lem:adelic_cech_d01_d12_formula} the
  degree-\(1\) differential \(d^1\), evaluated in each of the eight
  index components of \(C^2 = \prod_x \Gamma(\bigsqcap_a
  U_{x(a)}, F)\), collapses by the coface-restriction identities of
  \ref{mk-lem:adelic_cech_coface_restriction} (parallel restrictions
  cancel and diagonal restrictions round-trip). Consequently
  \(\ker d^1 \subseteq C^1 = \Gamma(U_0, F) \times \Gamma(U_1, F) \times
  \Gamma(U_1, F) \times \Gamma(U_0, F)\) (indexed by the four functions
  \(\mathrm{Fin}\,2 \to \{0,1\}\)) consists exactly of the tuples
  \((0, q, -q, 0)\) for \(q \in \Gamma(U_0 \cap U_1, F)\), giving a
  \(k\)-linear isomorphism \(\Gamma(U_0 \cap U_1, F) \xrightarrow{\ \sim\
  } \ker d^1\), \(q \mapsto (0, q, -q, 0)\). Under this identification
  the image of \(d^0 \colon \Gamma(U_0, F) \oplus \Gamma(U_1, F) \to
  \ker d^1\), \((f_0, f_1) \mapsto (0,\, f_0|_{01} - f_1|_{01},\,
  f_1|_{01} - f_0|_{01},\, 0)\), corresponds to \(\{ f_0|_{01} -
  f_1|_{01} : f_i \in \Gamma(U_i, F)\} = \Gamma(U_0, F) + \Gamma(U_1,
  F)\) inside \(\Gamma(U_0 \cap U_1, F)\) (each \(\Gamma(U_i, F)\) being
  closed under negation). Composing the two identifications gives the
  claimed \(k\)-linear isomorphism
  \(\mathrm{cechCohomology}\, C\, F\, \mathcal U\, 1 \simeq_k
  \Gamma(U_0 \cap U_1, F)/(\Gamma(U_0, F) + \Gamma(U_1, F))\).
\end{proof}

\begin{proposition}
[First cohomology vanishes on the affine pieces]
  \label{mk-thm:adelic_affine_h1_vanishing}
  \lean{AlgebraicGeometry.Scheme.subsingleton_hModule'_one_toModuleKSheaf_of_isAffineOpen}
  \leanok
  \dcref{custom}

  If \(U\subseteq C\) is affine, then
  \(H^1(U,\struct C|_U)=0\).  In particular this holds for both members
  of an affine Mayer--Vietoris square.
\end{proposition}
\begin{proof}
  \leanok
  This is affine vanishing for the quasi-coherent structure sheaf.
\end{proof}

\begin{proposition}
[The degree-one Mayer--Vietoris quotient]
  \label{mk-thm:adelic_mv_h1_cokernel_equiv}
  \lean{AlgebraicGeometry.Scheme.AffineCoverMVSquare.h1CokEquivHModule'One}
  \leanok
  \dcref{custom}

  \uses{mk-thm:Scheme_AffineCoverMVSquare_HModule_prime_sequence_exact}
  Let \(F\) be a sheaf of \(k\)-modules.  If
  \(H^1(U_0,F)=H^1(U_1,F)=0\), then the connecting homomorphism in the
  Mayer--Vietoris sequence induces an isomorphism
  \[
    \Gamma(U_0\cap U_1,F)/
      \operatorname{im}\bigl(\Gamma(U_0,F)\oplus\Gamma(U_1,F)\bigr)
      \simeq_k H^1(U_0\cup U_1,F).
  \]
\end{proposition}
\begin{proof}
  \leanok
  Exactness identifies the kernel of the connecting map with the image of
  the difference of restrictions.  Vanishing on the two pieces makes the
  connecting map surjective, so the first isomorphism theorem gives the
  displayed quotient.
\end{proof}

\begin{theorem}
[The structure-sheaf \(H^1\) is the cover cokernel]
  \label{mk-thm:adelic_HModuleOne_linearEquiv_cokernel}
  \lean{AlgebraicGeometry.Scheme.AffineCoverMVSquare.hModuleOneEquivH1Cok_curve}
  \leanok
  \dcref{custom}

  \uses{mk-thm:adelic_affine_h1_vanishing,
    mk-thm:adelic_mv_h1_cokernel_equiv}
  Every affine Mayer--Vietoris square on \(C\) gives a \(k\)-linear
  isomorphism
  \[
    \HModule\, k\, (\toModuleKSheaf C)\, 1
    \;\simeq_k\; \check{H}^1 \;=\; A_{01}/(A_0 + A_1).
  \]
\end{theorem}
\begin{proof}
  \leanok
  Apply \ref{mk-thm:adelic_mv_h1_cokernel_equiv} to the structure sheaf;
  its two vanishing hypotheses are \ref{mk-thm:adelic_affine_h1_vanishing}.
  Since \(U_0\cup U_1=C\), its target is the global first cohomology.
\end{proof}

\section{Finiteness via the projective line}
\label{mk-sec:RiemannRoch_Adelic_finiteness}

By \ref{mk-thm:adelic_HModuleOne_linearEquiv_cokernel}, it remains to prove
that \(A_{01}/(A_0+A_1)\) is finite-dimensional.  We pull back the two
standard charts of \(\mathbb P^1\) along a finite morphism.

\begin{proposition}
[A nonconstant map to \(\mathbb{P}^1\) is finite]
  \label{mk-thm:adelic_exists_finiteMorphismToP1}
  \lean{AlgebraicGeometry.Adelic.isFinite_left_of_exists_ne}
  \leanok
  \dcref{custom}
  \uses{mk-def:adelic_nonconstant_map_gate, mk-thm:adelic_nonconstant_map_isFinite}
  If a \(k\)-morphism \(\pi\colon C\to\mathbb P^1_k\) takes two distinct
  values, then \(\pi\) is finite.
\end{proposition}
\begin{proof}
  \leanok
  Let \(\pi \colon C \to \mathbb{P}^1_k\) be a nonconstant \(k\)-morphism
  (\ref{mk-def:adelic_nonconstant_map_gate}), with \(\pi(x_1) \ne
  \pi(x_2)\) for some \(x_1, x_2 \in C\). Since \(\mathbb{P}^1_k\) is
  separated and of finite type over \(k\), \ref{mk-thm:adelic_nonconstant_map_isFinite}
  applied to \(\pi\) shows \(\pi\) is a finite morphism.
\end{proof}

\begin{definition}
[A nonconstant map to \(\mathbb{P}^1\)]
  \label{mk-def:adelic_nonconstant_map_gate}
  \lean{AlgebraicGeometry.Adelic.ExistsNonconstantMapToP1}
  \leanok
  \dcref{custom}

  \uses{mk-def:projective_space}
  The class \(\mathrm{ExistsNonconstantMapToP1}\, C\) asserts the
  existence of a \(k\)-morphism \(\pi \colon C \to \mathbb{P}^1_k\)
  taking two distinct values.
\end{definition}

\begin{proposition}
[A nonconstant map to the integral projective line]
  \label{mk-thm:adelic_exists_nonconstant_projint}
  \lean{AlgebraicGeometry.Adelic.ExistsNonconstantMapToProjInt,
    AlgebraicGeometry.Adelic.existsNonconstantMapToProjInt_of_ajc}
  \leanok
  \dcref{custom}

  A smooth proper geometrically integral curve admits a morphism
  \(f\colon C\to\operatorname{Proj}\mathbb Z[X_0,X_1]\) taking two
  distinct values.
\end{proposition}
\begin{proof}
  \leanok
  A transcendental element of \(K/k\) defines a rational map to the
  projective line.  Properness and regularity in dimension one extend it
  across every closed point, and transcendence makes the map nonconstant.
\end{proof}

\begin{proposition}
[Integral-model bridge to \(\mathbb P^1_k\)]
  \label{mk-lem:adelic_nonconstant_projint_bridge}
  \lean{AlgebraicGeometry.Adelic.existsNonconstantMapToP1_of_nonconstantMapToProjInt}
  \leanok
  \dcref{custom}
  \uses{mk-thm:adelic_exists_nonconstant_projint, mk-def:adelic_nonconstant_map_gate}
  \(\mathrm{ExistsNonconstantMapToProjInt}\, C\) implies
  \(\mathrm{ExistsNonconstantMapToP1}\, C\).
\end{proposition}
\begin{proof}
  \leanok
  The concrete \(\mathbb{P}^1_k = \mathbb{P}(\mathrm{ULift}(\mathrm{Fin}\,
  2); \operatorname{Spec} k)\) is, by construction, the fibre product
  \(\operatorname{Spec} k \times_{\top} \operatorname{Proj}
  \mathbb{Z}[X_0, X_1]\) over the \emph{terminal} scheme \(\top\). Since
  any two morphisms into \(\top\) agree, the compatibility square
  required by the universal property of this pullback is automatic:
  given \(f \colon C \to \operatorname{Proj} \mathbb{Z}[X_0, X_1]\) and
  the ambient structure map \(C \to \operatorname{Spec} k\), the two
  assemble into a unique \(g \colon C \to \mathbb{P}^1_k\) with
  \(g\) over \(\operatorname{Spec} k\) and \(g\) composed with the
  projection \(\mathbb{P}^1_k \to \operatorname{Proj} \mathbb{Z}[X_0,
  X_1]\) equal to \(f\). If \(f(x_1) \ne f(x_2)\) then, composing with
  the projection, \(g(x_1) \ne g(x_2)\); hence \(g\) is a nonconstant
  \(k\)-morphism \(C \to \mathbb{P}^1_k\).
\end{proof}

\begin{lemma}
[Quasi-finite algebras do not raise Krull dimension]
  \label{mk-lem:adelic_quasiFinite_dim_le}
  \lean{AlgebraicGeometry.Adelic.ringKrullDim_le_of_quasiFinite}
  \leanok
  \dcref{custom}
  For a quasi-finite \(R\)-algebra \(A\), \(\operatorname{ringKrullDim}
  A \le \operatorname{ringKrullDim} R\).
\end{lemma}
\begin{proof}
  \leanok
  Quasi-finiteness gives incomparability of primes over the contraction
  map \(\operatorname{Spec} A \to \operatorname{Spec} R\): if two primes
  of \(A\) contract to the same prime of \(R\) and one is contained in
  the other, they are equal. Hence the contraction map is strictly
  monotone on chains of primes, so it strictly increases no chain
  length in the wrong direction and \(\operatorname{ringKrullDim} A \le
  \operatorname{ringKrullDim} R\).
\end{proof}

\begin{lemma}
[Dimension one for standard-smooth curves]
  \label{mk-lem:adelic_relDimOne_dim_le_one}
  \lean{AlgebraicGeometry.Adelic.ringKrullDim_le_one_of_isStandardSmoothOfRelativeDimension_one}
  \leanok
  \dcref{custom}
  \uses{mk-lem:adelic_quasiFinite_dim_le}
  Let \(S\) be standard-smooth of relative dimension \(1\) over a field
  \(k\), and let \(q\) be a prime of \(S\). Every localisation \(S_q\)
  of \(S\) at \(q\) satisfies \(\operatorname{ringKrullDim} S_q \le 1\).
\end{lemma}
\begin{proof}
  \leanok
  Standard smoothness of relative dimension one makes \(\Omega_{S/k}\)
  free of rank one; writing the spanning vector as a finite combination
  of exact differentials, some \(dt\) (for a generator \(t\) of \(S\))
  has coefficient \(a \notin q\) in this basis. Over \(S_q\) the element
  \(a\) is a unit, so \(dt\) spans \(\Omega_{S_q/k}\); the \(k\)-algebra
  map \(k[X] \to S_q\), \(X \mapsto t\), then makes \(S_q\) formally
  unramified over \(k[X]\) (\(\Omega_{S_q/k[X]} = 0\)) and essentially
  of finite type, i.e.\ \emph{quasi-finite} over \(k[X]\). Since
  \(k[X]\) is a PID, \(\operatorname{ringKrullDim} k[X] = 1\), and
  \ref{mk-lem:adelic_quasiFinite_dim_le} gives
  \(\operatorname{ringKrullDim} S_q \le 1\).
\end{proof}

\begin{lemma}
[Stalks of the curve have dimension \(\le 1\)]
  \label{mk-lem:adelic_curve_stalk_dim_le_one}
  \lean{AlgebraicGeometry.Adelic.ringKrullDim_stalk_le_one_of_curve}
  \leanok
  \dcref{custom}
  \uses{mk-lem:adelic_relDimOne_dim_le_one}
  Every stalk \(\struct{C, z}\) of the curve \(C\) (smooth of relative
  dimension one over \(\operatorname{Spec} k\)) has
  \(\operatorname{ringKrullDim} \struct{C, z} \le 1\).
\end{lemma}
\begin{proof}
  \leanok
  As \(\operatorname{Spec} k\) is a single point, \(z\) has a standard
  smooth chart \(V\) of relative dimension one over all of
  \(\operatorname{Spec} k\); the stalk \(\struct{C, z}\) is the
  localisation of \(\Gamma(V, \struct{C})\) at the prime of \(z\), so
  \ref{mk-lem:adelic_relDimOne_dim_le_one} applies directly.
\end{proof}

\begin{lemma}
[Coheight \(\le 1\) on the curve]
  \label{mk-lem:adelic_curve_coheight_le_one}
  \lean{AlgebraicGeometry.Adelic.coheight_le_one_of_curve}
  \leanok
  \dcref{custom}
  \uses{mk-lem:adelic_curve_stalk_dim_le_one, mk-thm:ringKrullDim_stalk_eq_coheight}
  Every point \(z\) of \(C\) has topological coheight
  \(\operatorname{coheight} z \le 1\).
\end{lemma}
\begin{proof}
  \leanok
  Immediate from \ref{mk-lem:adelic_curve_stalk_dim_le_one} and the
  stalk--coheight bridge \ref{mk-thm:ringKrullDim_stalk_eq_coheight}:
  \[
    \operatorname{coheight}z
      =\operatorname{ringKrullDim}\struct{C,z}.
  \]
\end{proof}

\begin{lemma}
[Non-generic points of the curve are closed]
  \label{mk-lem:adelic_nongeneric_closed}
  \lean{AlgebraicGeometry.Adelic.isClosed_singleton_of_coheight_le_one}
  \leanok
  \dcref{custom}
  \uses{mk-lem:adelic_curve_coheight_le_one}
  On an irreducible scheme all of whose points have coheight \(\le 1\),
  every point \(z\) other than the generic point is closed.
\end{lemma}
\begin{proof}
  \leanok
  If \(\{z\}\) were not closed, its closure would contain a point
  \(z' \ne z\) with \(z \rightsquigarrow z'\), giving a strict
  specialisation chain \(z' \prec z \prec \eta\) (\(\eta\) the generic
  point, since \(z \ne \eta\)); each strict step adds at least one to
  coheight, so \(\operatorname{coheight} z' \ge 2\), contradicting the
  hypothesis.
\end{proof}

\begin{lemma}
[Proper closed subsets of the curve are finite]
  \label{mk-lem:adelic_proper_closed_finite}
  \lean{AlgebraicGeometry.Adelic.finite_of_isClosed_of_ne_univ}
  \leanok
  \dcref{custom}
  \uses{mk-lem:adelic_nongeneric_closed}
  On an irreducible Noetherian sober scheme all of whose points have
  coheight \(\le 1\), every closed subset \(F \ne C\) is finite.
\end{lemma}
\begin{proof}
  \leanok
  A Noetherian space has finitely many irreducible components; a proper
  closed \(F\) does not contain the generic point, so each of its
  (finitely many) irreducible components has a generic point \(\zeta
  \ne \eta\), which is closed by \ref{mk-lem:adelic_nongeneric_closed}.
  A sober irreducible closed set with closed generic point is that
  singleton, so each component is a single point and \(F\) is finite.
\end{proof}

\begin{lemma}
[Fibres of a nonconstant map are finite]
  \label{mk-lem:adelic_nonconstant_finite_fibres}
  \lean{AlgebraicGeometry.Adelic.finite_preimage_singleton_of_exists_ne}
  \leanok
  \dcref{custom}
  \uses{mk-lem:adelic_proper_closed_finite}
  Let \(\pi \colon C \to Y\) be a \(k\)-morphism to a
  locally-of-finite-type \(k\)-scheme \(Y\) taking two distinct values.
  Every fibre \(\pi^{-1}(y)\) is finite.
\end{lemma}
\begin{proof}
  \leanok
  If \(\{y\}\) is closed, \(\pi^{-1}(y)\) is a closed subset of \(C\)
  that is not all of \(C\) (since \(\pi\) is nonconstant), hence finite
  by \ref{mk-lem:adelic_proper_closed_finite}. If \(\{y\}\) is not closed,
  then \(y\) is not a closed point of the Jacobson scheme \(Y\); since
  closed points of \(C\) map to closed points of \(Y\), the fibre
  \(\pi^{-1}(y)\) avoids every closed point of \(C\)
  (\ref{mk-lem:adelic_nongeneric_closed}: every non-generic point of
  \(C\) is closed), so \(\pi^{-1}(y) \subseteq \{\eta\}\) is finite.
\end{proof}

\begin{proposition}
[Nonconstant morphisms from the curve are finite]
  \label{mk-thm:adelic_nonconstant_map_isFinite}
  \lean{AlgebraicGeometry.Adelic.isFinite_left_of_exists_ne}
  \leanok
  \dcref{custom}
  \uses{mk-lem:adelic_nonconstant_finite_fibres}
  Let \(Y\) be a separated, locally-of-finite-type \(k\)-scheme. Any
  \(k\)-morphism \(\pi \colon C \to Y\) taking two distinct values is a
  \emph{finite} morphism.
\end{proposition}
\begin{proof}
  \leanok
  \(\pi\) is proper: \(\pi\) composed with the separated structure map
  \(Y \to \operatorname{Spec} k\) is the proper structure map
  \(C \to \operatorname{Spec} k\), and properness cancels along a
  separated morphism, so \(\pi\) itself is proper.
  \(\pi\) is locally quasi-finite by finiteness of its fibres
  (\ref{mk-lem:adelic_nonconstant_finite_fibres}). A proper, locally
  quasi-finite morphism is finite (Zariski's main theorem), so \(\pi\)
  is finite.
\end{proof}

\begin{proposition}
[Charts are module-finite over the \(\mathbb{P}^1\) charts]
  \label{mk-thm:adelic_chart_finite_over_P1}
  \lean{AlgebraicGeometry.Adelic.isAffineOpen_preimage_of_isFinite,
    AlgebraicGeometry.Adelic.module_finite_app_of_isFinite,
    AlgebraicGeometry.Adelic.LaurentChartData.pullbackSquare}
  \leanok
  \dcref{custom}
  \uses{mk-thm:adelic_exists_finiteMorphismToP1, mk-thm:adelic_chartRing_isDedekindDomain}
  Take \(x\) as in \ref{mk-thm:adelic_exists_finiteMorphismToP1} and the
  cover \(U_0 = \pi^{-1}(\mathbb{A}^1_x)\),
  \(U_1 = \pi^{-1}(\mathbb{A}^1_{1/x})\); both are affine, since the
  preimage of an affine open under a finite (hence affine) morphism is
  affine. Then \(A_0 = \Gamma(U_0, \struct{C})\) is a finite module over
  \(k[x]\), and \(A_1 = \Gamma(U_1, \struct{C})\) is a finite module over
  \(k[1/x]\): for a finite morphism \(\rho\) and an affine open \(V\) of
  its target \(Y\), \(\Gamma(\rho^{-1}V, \struct{X})\) is always a finite
  module over \(\Gamma(V, \struct{Y})\) (the integral-closure finiteness
  of a finite morphism, applied to \(\rho = \pi\), \(Y = \mathbb{P}^1_k\)).
  Packaging \(U_0, U_1\) and their affine overlap as a
  \(2\)-affine cover of \(C\) is the pulled-back cover square consumed by
  the keystone (\ref{mk-thm:adelic_finiteDimensional_cokernel_zero}). (The
  sharper classical description of \(A_0\), \(A_1\), and \(A_{01}\) as
  integral closures, generically free of rank \(n=[K:k(x)]\), is not
  needed for this argument.)
\end{proposition}
\begin{proof}
  \leanok
  Affineness of \(U_0\) and \(U_1\) is preimage-stability of affine opens
  under a finite morphism. Module-finiteness of \(A_0\) over \(k[x]\)
  (and symmetrically of \(A_1\) over \(k[1/x]\)) is the general fact that
  for a finite morphism \(\rho \colon X \to Y\) and an affine open \(V\)
  of \(Y\), the induced ring map \(\Gamma(V, \struct{Y}) \to
  \Gamma(\rho^{-1}V, \struct{X})\) makes the target a finite module over
  the source.
\end{proof}
\begin{lemma}
[The \(\mathbb{P}^1\) base cokernel vanishes]
  \label{mk-thm:adelic_cokernel_P1_zero}
  \lean{AlgebraicGeometry.Adelic.nonneg_sup_nonpos_eq_top,
    AlgebraicGeometry.Adelic.subsingleton_laurentCokernel}
  \leanok
  \dcref{custom}

  The untwisted cokernel of the standard cover of \(\mathbb{P}^1_k\)
  vanishes:
  \[
    k[x, x^{-1}] \big/ \bigl(k[x] + k[x^{-1}]\bigr) \;=\; 0,
    \qquad\text{i.e.}\qquad H^1(\mathbb{P}^1_k, \struct{\mathbb{P}^1}) = 0.
  \]
\end{lemma}
\begin{proof}
  \leanok
  As a \(k\)-vector space \(k[x, x^{-1}]\) has basis the Laurent
  monomials \(\{x^m : m \in \mathbb{Z}\}\). Every monomial \(x^m\) with
  \(m \ge 0\) lies in \(k[x]\), and every monomial with \(m \le 0\) lies
  in \(k[x^{-1}]\); since each integer is \(\ge 0\) or \(\le 0\), the sum
  \(k[x] + k[x^{-1}]\) contains every basis monomial and hence equals
  \(k[x, x^{-1}]\). Therefore the quotient is zero. This is the genus-zero
  statement for \(\mathbb{P}^1_k\).
\end{proof}

\begin{theorem}
[The two-cover cokernel of the curve is finite-dimensional]
  \label{mk-thm:adelic_finiteDimensional_cokernel_zero}
  \lean{AlgebraicGeometry.Adelic.LaurentChartData,
    AlgebraicGeometry.Adelic.LaurentChartData.pullbackSquare,
    AlgebraicGeometry.Adelic.LaurentChartData.module_finite_H1Cok,
    AlgebraicGeometry.Adelic.exists_affineCoverMVSquare_module_finite_H1Cok}
  \leanok
  \dcref{custom}

  \uses{mk-thm:adelic_exists_finiteMorphismToP1, mk-thm:adelic_chart_finite_over_P1,
    mk-thm:adelic_cokernel_P1_zero, mk-def:adelic_coverH1}
  Let \(\pi \colon C \to Y\) be a finite \(k\)-morphism of
  \(\operatorname{Spec} k\)-schemes whose target carries \emph{Laurent
  chart data}: affine opens \(V_0 \cup V_1 = Y\), coordinates
  \(x \in \Gamma(V_0)\), \(y \in \Gamma(V_1)\) with
  \(V_0 \cap V_1 = D(x) = D(y)\), \(x \cdot y = 1\) on the overlap, and
  \(\Gamma(V_0) = \sum_n k\, x^n\), \(\Gamma(V_1) = \sum_n k\, y^n\) (for
  \(Y = \mathbb{P}^1_k\): \(\Gamma(V_0) = k[x]\),
  \(\Gamma(V_1) = k[x^{-1}]\)). Set \(U_i = \pi^{-1} V_i\),
  \(A_i = \Gamma(U_i, \struct{C})\),
  \(M = \Gamma(U_0 \cap U_1, \struct{C})\). Then the two-cover Čech
  cokernel \(\check{H}^1 = M/(A_0 + A_1)\) is a finite-dimensional
  \(k\)-vector space. In particular, under the gates
  \(\mathrm{HasFiniteMapToP1}\, C\) and \(\mathrm{P1HasLaurentChartData}\,
  k\), some \(2\)-affine cover of \(C\) (the pullback of the standard
  charts) has finite-dimensional \(\check{H}^1\).
\end{theorem}
\begin{proof}
  \leanok
  Write \(t = \pi^{\sharp} x|_{U_0 \cap U_1}\),
  \(u = \pi^{\sharp} y|_{U_0 \cap U_1}\); by naturality of \(\pi^{\sharp}\)
  and \(x \cdot y = 1\) they satisfy \(t \cdot u = 1\) in \(M\) (an
  abstract Laurent pair). Let \(N_0, N_1 \subseteq M\) be the images of
  the restrictions \(A_0 \to M\), \(A_1 \to M\); \(N_0\) is stable under
  multiplication by \(t\) and \(N_1\) under multiplication by \(u\),
  since restriction is a ring map.

  \emph{Extension.} \(U_0 \cap U_1 = \pi^{-1}(V_0 \cap V_1) =
  \pi^{-1}(D(x)) = D(\pi^{\sharp} x)\) is the basic open of
  \(\pi^{\sharp} x\) inside the affine \(U_0\), so \(M\) is the
  localization \((A_0)_t\); hence for every \(m \in M\) there are \(n\)
  and \(a \in A_0\) with \(t^n m = a|_{U_0 \cap U_1}\), and symmetrically
  \(u^n m \in N_1\) for \(n \gg 0\).

  \emph{Ladder spanning.} By finiteness of \(\pi\), \(A_0\) is a finite
  \(\Gamma(V_0)\)-module; choose generators \(g_1, \dots, g_r\). Since
  \(\Gamma(V_0)\) is the \(k\)-span of the powers of \(x\) (the geometric
  source of the Laurent split of \ref{mk-thm:adelic_cokernel_P1_zero}),
  \(A_0\) is the \(k\)-span of \(\{x_0^j g_i\}\) where
  \(x_0 = \pi^{\sharp} x\); combining with the extension property and
  \(tu = 1\), every \(m \in M\) lies in the \(k\)-span of the two-sided
  ladder \(\{t^j \bar{g}_i\} \cup \{u^j \bar{g}_i\}\), \(j \in
  \mathbb{N}\), with \(\bar{g}_i \in N_0\) the images of the generators.

  \emph{Middle band.} The rungs \(t^j \bar{g}_i\) lie in \(N_0\)
  (\(t\)-stability); choosing \(n_i\) with \(u^{n_i} \bar{g}_i \in N_1\)
  and \(N = \max_i n_i\), all rungs \(u^j \bar{g}_i\) with \(j \ge N\) lie
  in \(N_1\) (\(u\)-stability). Hence modulo \(N_0 + N_1\) the module
  \(M\) is spanned by the finitely many middle rungs
  \(\{u^j \bar{g}_i : i \le r,\ j < N\}\), so
  \(\check{H}^1 = M/(N_0 + N_1)\) is finite-dimensional over \(k\).
\end{proof}

\subsection{The projective-line chart rings}
\label{mk-sec:RiemannRoch_Adelic_P1ChartData}

The two standard charts of \(\mathbb P^1_k\) have coordinate rings
\(k[x]\) and \(k[x^{-1}]\), and their intersection has coordinate ring
\(k[x,x^{-1}]\).  The following construction obtains these charts by
base change from the degree-zero away rings of
\(\operatorname{Proj}\mathbb Z[X_0,X_1]\).

\begin{definition}
[The coordinate fraction \(X_j/X_i\)]
  \label{mk-def:adelic_p1CoordAway}
  \lean{AlgebraicGeometry.Adelic.p1CoordAway}
  \leanok
  \dcref{custom}
  \uses{mk-def:projective_space}
  For \(i, j \in \{0, 1\}\), the coordinate fraction \(X_j/X_i\) is the
  element of the degree-zero away ring \((\mathbb{Z}[X_0,X_1]_{X_i})_0\)
  with numerator \(X_j\) and denominator \(X_i\) (both homogeneous of
  degree \(1\)).
\end{definition}

\begin{definition}
[The pulled-back chart coordinates]
  \label{mk-def:adelic_p1_chart_sections}
  \lean{AlgebraicGeometry.Adelic.p1XSection, AlgebraicGeometry.Adelic.p1YSection}
  \leanok
  \dcref{custom}
  \uses{mk-def:adelic_p1CoordAway}
  The chart coordinates \(x = X_1/X_0 \in \Gamma(V_0)\) and
  \(y = X_0/X_1 \in \Gamma(V_1)\) on the standard two-chart cover
  \(V_0, V_1\) of the concrete model \(\mathbb{P}^1_k\), obtained by
  pulling back the away-section fractions
  \(X_1/X_0 \in \mathrm{Away}\,\mathcal A\,X_0\) and
  \(X_0/X_1 \in \mathrm{Away}\,\mathcal A\,X_1\)
  (\ref{mk-def:adelic_p1CoordAway}) along the projection
  \(\mathbb{P}^1_k \to \operatorname{Proj}\mathbb{Z}[X_0,X_1]\).
\end{definition}

\begin{lemma}
[The monomial-to-power identity]
  \label{mk-lem:adelic_p1Monomial_eq_coordPow}
  \lean{AlgebraicGeometry.Adelic.p1Monomial_eq_coordPow}
  \leanok
  \dcref{custom}
  \uses{mk-def:adelic_p1CoordAway}
  For \(a, b \in \mathbb{N}\), the degree-\((a+b)\) monomial fraction
  \(X_0^a X_1^b / X_0^{a+b}\) in \((\mathbb{Z}[X_0,X_1]_{X_0})_0\) equals
  \((X_1/X_0)^b\).
\end{lemma}
\begin{proof}
  \leanok
  Both sides are the class of \(X_0^a X_1^b\) over the denominator
  \(X_0^{a+b}\) in the localisation: cross-multiplying,
  \((X_0)^b \cdot (X_0^a X_1^b) = X_0^{a+b} \cdot X_1^b\), a monomial
  identity in the commutative monoid of monomials.
\end{proof}

\begin{theorem}
[Span of the projective-line chart]
  \label{mk-thm:adelic_span_p1CoordAway_pow_top}
  \lean{AlgebraicGeometry.Adelic.span_p1CoordAway_pow_top}
  \leanok
  \dcref{custom}
  \uses{mk-lem:adelic_p1Monomial_eq_coordPow}
  The degree-zero away ring \((\mathbb{Z}[X_0,X_1]_{X_0})_0\) is spanned,
  as a module over its constant subring \((\mathbb{Z}[X_0,X_1])_0\), by
  the powers of the coordinate fraction \(X_1/X_0\).
\end{theorem}
\begin{proof}
  \leanok
  Every element of the away ring is a fraction \(p/X_0^N\) with \(p\)
  homogeneous of degree \(N\) (\(N \ge 0\)). Expanding \(p\) over its
  monomial support \(\{c_d X_0^{d_0} X_1^{d_1}\}\), homogeneity forces
  \(d_0 + d_1 = N\) for every occurring exponent vector \(d\), so each
  summand is, by \ref{mk-lem:adelic_p1Monomial_eq_coordPow}, the constant
  \(c_d\) times \((X_1/X_0)^{d_1}\). Summing over the (finite) support
  writes \(p/X_0^N\) as a \((\mathbb{Z}[X_0,X_1])_0\)-combination of
  powers of \(X_1/X_0\).
\end{proof}

\begin{theorem}
[Laurent chart data on \(\mathbb P^1\)]
  \label{mk-thm:adelic_P1_laurent_chart_data}
  \lean{AlgebraicGeometry.Adelic.P1HasLaurentChartData,
    AlgebraicGeometry.Adelic.p1CoverSquare,
    AlgebraicGeometry.Adelic.p1LaurentChartData,
    AlgebraicGeometry.Adelic.instP1HasLaurentChartData}
  \leanok
  \dcref{custom}
  \uses{mk-def:projective_space, mk-thm:adelic_span_p1CoordAway_pow_top}
  For every field \(k\), the standard two-chart cover of
  \(\mathbb P^1_k\) carries Laurent chart data: its chart rings are
  spanned by nonnegative powers of \(x\) and \(x^{-1}\), the overlap is
  both \(D(x)\) and \(D(x^{-1})\), and the two coordinates multiply to
  one there.
\end{theorem}

\begin{theorem}
[Finite-dimensionality of \(H^1(C,\struct C)\)]
  \label{mk-thm:adelic_module_finite_genus}
  \lean{AlgebraicGeometry.Adelic.module_finite_hModule_one,
    AlgebraicGeometry.Adelic.instModuleFiniteHModuleOne}
  \leanok
  \dcref{custom}
  \uses{mk-thm:adelic_HModuleOne_linearEquiv_cokernel,
    mk-thm:adelic_finiteDimensional_cokernel_zero,
    mk-thm:adelic_exists_finiteMorphismToP1,
    mk-thm:adelic_P1_laurent_chart_data, mk-def:genus}
  If \(C\) is smooth and proper with geometrically integral fibres, then
  \[
    \Module.\Finite\ k\ \bigl(\HModule\, k\, (\toModuleKSheaf C)\, 1\bigr).
  \]
  Thus \(\genus C\) is the dimension of this finite-dimensional vector
  space.
\end{theorem}
\begin{proof}
  \leanok
  Pull back the standard charts along the finite morphism to
  \(\mathbb P^1_k\).  Their cokernel is finite-dimensional by
  \ref{mk-thm:adelic_finiteDimensional_cokernel_zero}, and
  \ref{mk-thm:adelic_HModuleOne_linearEquiv_cokernel} transports this
  finiteness to \(H^1(C,\struct C)\).
\end{proof}

\section{The Euler-characteristic ledger}
\label{mk-sec:RiemannRoch_Adelic_chi}

Fix the same two-affine cover.  The divisor section groups become
\(k\)-subspaces of \(K\), and their quotients give a concrete numerical
Euler-characteristic calculation.

\subsection{The twist sequence}

\begin{lemma}
[Antitonicity of the section subgroup in the open]
  \label{mk-lem:adelic_sectionOfDivisor_antitone_open}
  \lean{AlgebraicGeometry.Adelic.sectionOfDivisor_antitone_open}
  \leanok
  \dcref{custom}
  \uses{mk-def:adelic_sectionOfDivisor}
  For opens \(U \subseteq U'\) of \(C\) and a Weil divisor \(D\),
  \(\Gamma(U', \struct{C}(D)) \subseteq \Gamma(U, \struct{C}(D))\): a
  larger open imposes the order bound \(v_P(f) \ge -D(P)\) at more
  prime divisors, so its section subgroup is smaller.
\end{lemma}
\begin{proof}
  \leanok
  A nonzero \(f \in \Gamma(U', \struct{C}(D))\) satisfies
  \(v_P(f) \ge -D(P)\) at every prime divisor \(P\) with point in
  \(U'\), in particular at every such \(P\) with point in the smaller
  open \(U \subseteq U'\), so \(f \in \Gamma(U, \struct{C}(D))\).
\end{proof}

\begin{definition}
[The coboundary subgroup]
  \label{mk-def:adelic_coboundary}
  \lean{AlgebraicGeometry.Adelic.coboundary}
  \leanok
  \dcref{custom}
  \uses{mk-def:adelic_sectionOfDivisor}
  The \emph{coboundary} is the sum of the two chart section subgroups
  inside \(K\):
  \[
    B(D)=\Gamma(U_0,\struct C(D))+\Gamma(U_1,\struct C(D)).
  \]
\end{definition}

\begin{lemma}
[The coboundary and the linear system land in the overlap]
  \label{mk-lem:adelic_coboundary_le_overlap}
  \lean{AlgebraicGeometry.Adelic.coboundary_le_overlap,
    AlgebraicGeometry.Adelic.linearSystem_le_overlap}
  \leanok
  \dcref{custom}
  \uses{mk-def:adelic_coboundary, mk-def:adelic_sectionOfDivisor,
    mk-lem:adelic_sectionOfDivisor_antitone_open}
  \(B(D) \subseteq \Gamma(U_0 \cap U_1, \struct{C}(D))\) and
  \(L(D) \subseteq \Gamma(U_0 \cap U_1, \struct{C}(D))\): each of the two
  chart sections, and every global section, restricts to the overlap.
\end{lemma}
\begin{proof}
  \leanok
  Immediate from the antitonicity of \(\Gamma(-, \struct{C}(D))\) in the
  open (\ref{mk-lem:adelic_sectionOfDivisor_antitone_open}: a larger open
  imposes more order conditions, hence restricts into a smaller open's
  section group), applied to \(U_0 \cap U_1 \subseteq U_0\), \(U_0 \cap
  U_1 \subseteq U_1\), and \(U_0 \cap U_1 \subseteq \top\).
\end{proof}

\begin{theorem}
[Exactness at the \(L\)-terms]
  \label{mk-thm:adelic_linearSystem_inf_overlap_eq}
  \lean{AlgebraicGeometry.Adelic.linearSystem_inf_overlap_eq}
  \leanok
  \dcref{custom}
  \uses{mk-lem:adelic_coboundary_le_overlap}
  Suppose \(D \le D'\) pointwise and \(D' - D\) is supported on the
  overlap \(V = U_0 \cap U_1\) (off \(V\), \(D\) and \(D'\) agree). Then
  \[
    L(D') \cap \Gamma(V, \struct{C}(D)) \;=\; L(D).
  \]
\end{theorem}
\begin{proof}
  \leanok
  (\(\supseteq\)) A global section of \(\struct{C}(D)\) is a global
  section of \(\struct{C}(D')\) (\(D \le D'\)) that trivially restricts
  into \(\Gamma(V, \struct{C}(D))\). (\(\subseteq\)) Let \(f \in L(D')\)
  restrict into \(\Gamma(V, \struct{C}(D))\); if \(f = 0\) it lies in
  \(L(D)\) trivially. Otherwise, to check \(f \in L(D)\) it suffices to
  bound \(v_P(f) \ge -D(P)\) at every prime divisor \(P\): on \(V\) this
  is the hypothesis \(f \in \Gamma(V, \struct{C}(D))\), and off \(V\) it
  is the bound \(v_P(f) \ge -D'(P) = -D(P)\) (using \(D = D'\) off
  \(V\)) already known from \(f \in L(D')\).
\end{proof}

\begin{definition}
[The function-field quotient and the window map]
  \label{mk-def:adelic_H1_window}
  \lean{AlgebraicGeometry.Adelic.H1, AlgebraicGeometry.Adelic.H1Twist,
    AlgebraicGeometry.Adelic.windowMap}
  \leanok
  \dcref{custom}
  \uses{mk-def:adelic_function_field_H1, mk-def:adelic_coboundary,
    mk-lem:adelic_coboundary_le_overlap}
  The function-field quotient \(\check{H}^1_{\mathrm{ff}}(D) = \Gamma(V,
  \struct{C}(D))/B(D)\) (\(B(D) \subseteq \Gamma(V, \struct{C}(D))\) by
  \ref{mk-lem:adelic_coboundary_le_overlap}). For \(D \le D'\), the
  inclusions \(\Gamma(V, \struct{C}(D)) \subseteq \Gamma(V,
  \struct{C}(D'))\) and \(B(D) \subseteq B(D')\) induce the twist map
  \(\check{H}^1_{\mathrm{ff}}(D) \to
  \check{H}^1_{\mathrm{ff}}(D')\), and the inclusions \(L(D)
  \subseteq \Gamma(V, \struct{C}(D))\) induce the \emph{window map}
  \[
    L(D')/L(D) \longrightarrow \Gamma(V, \struct{C}(D'))\big/
      \Gamma(V, \struct{C}(D)),
  \]
  the left-hand map of the four-term ledger sequence of
  \ref{mk-thm:adelic_twist_exact_sequence}.
\end{definition}

\begin{theorem}
[The window map is injective]
  \label{mk-thm:adelic_windowMap_injective}
  \lean{AlgebraicGeometry.Adelic.windowMap_injective}
  \leanok
  \dcref{custom}
  \uses{mk-def:adelic_H1_window, mk-thm:adelic_linearSystem_inf_overlap_eq}
  Under the hypotheses of \ref{mk-thm:adelic_linearSystem_inf_overlap_eq}
  (\(D \le D'\), \(D' - D\) supported on the overlap), the window map
  \(L(D')/L(D) \to \Gamma(V, \struct{C}(D'))/\Gamma(V, \struct{C}(D))\)
  is injective: this is the exactness of the ledger sequence at the two
  \(L\)-terms.
\end{theorem}
\begin{proof}
  \leanok
  A class \([f] \in L(D')/L(D)\) mapping to zero means \(f \in L(D')\)
  restricts into \(\Gamma(V, \struct{C}(D))\); by
  \ref{mk-thm:adelic_linearSystem_inf_overlap_eq} this forces \(f \in
  L(D)\), i.e.\ \([f] = 0\).
\end{proof}

\begin{definition}[The divisorial sheaf]
  \label{mk-def:adelic_divisorial_sheaf}
  \dcref{custom}
  \uses{mk-def:adelic_sectionOfDivisor}
  The divisorial sheaf \(\struct C(D)\) is the subsheaf of the constant
  sheaf \(K\) whose sections on a nonempty open \(U\) are the rational
  functions satisfying
  \[
    v_P(f)+D(P)\geq 0
    \quad\text{for every prime divisor }P\text{ meeting }U.
  \]
  Its section over the empty open is zero.  It is an invertible, hence
  quasi-coherent, \(\struct C\)-module.
\end{definition}

\begin{proposition}[Comparison with the divisorial sheaf]
  \label{mk-thm:adelic_function_field_H1_compare}
  \dcref{custom}
  \uses{mk-def:adelic_divisorial_sheaf,
    mk-def:adelic_function_field_H1,
    mk-thm:adelic_mv_h1_cokernel_equiv}
  For every divisor \(D\), the order-bounded function-field quotient is
  canonically isomorphic to the first cohomology of the divisorial sheaf:
  \[
    \check H^1_{\mathrm{ff}}(D)\simeq_k H^1(C,\struct C(D)).
  \]
\end{proposition}
\begin{proof}
  The valuation description of \(\struct C(D)\) identifies its sections on
  \(U_0\), \(U_1\), and \(U_0\cap U_1\) with the three divisor section
  groups.  The resulting difference complex has cokernel
  \(\check H^1_{\mathrm{ff}}(D)\).  Since \(\struct C(D)\) is
  quasi-coherent and both charts are affine, Mayer--Vietoris identifies
  this cokernel with \(H^1(C,\struct C(D))\).
\end{proof}

\begin{proposition}
[Twist exact sequence]
  \label{mk-thm:adelic_twist_exact_sequence}
  \dcref{custom}
  \uses{mk-def:adelic_divisorial_sheaf,
    mk-thm:adelic_function_field_H1_compare}
  For Weil divisors \(D \le D'\) there is a four-term exact sequence of
  \(k\)-vector spaces
  \[
    0 \longrightarrow L(D')/L(D)
      \longrightarrow \mathcal{Q}(D, D')
      \longrightarrow \check{H}^1_{\mathrm{ff}}(D)
      \longrightarrow \check{H}^1_{\mathrm{ff}}(D')
      \longrightarrow 0,
  \]
  where \(\mathcal{Q}(D, D') = H^0\bigl(C, \struct{C}(D')/\struct{C}(D)\bigr)\)
  is the space of global sections of the finite-support quotient sheaf.
\end{proposition}
\begin{proof}
  The inclusion of divisorial sheaves has finite-support cokernel
  \(\mathcal T\):
  \[
    0\longrightarrow\struct C(D)\longrightarrow\struct C(D')
      \longrightarrow\mathcal T\longrightarrow0.
  \]
  A finite-support sheaf has no higher cohomology.  The associated long
  exact sequence therefore ends after \(H^1(C,\struct C(D'))\), and the
  degree-zero terms identify with \(L(D)\), \(L(D')\), and
  \(\mathcal Q(D,D')\).  Since the cover is affine and the divisorial
  sheaves are quasi-coherent, their first sheaf cohomology groups are the
  displayed cover cokernels.  Quotienting the first two terms gives the
  stated sequence.
\end{proof}

\subsection{The one-point step}

\begin{definition}
[The single-point order-\(\ge m\) subgroup]
  \label{mk-def:adelic_orderGe}
  \lean{AlgebraicGeometry.Adelic.orderGe}
  \leanok
  \dcref{custom}
  \uses{mk-def:adelic_pointValuation}
  For a prime divisor \(P\) and \(m \in \mathbb{Z}\),
  \(\mathfrak G_P^{\ge m} = \{ f \in K : f = 0 \text{ or } v_P(f) \ge m
  \}\), an additive subgroup of \(K\) (for \(m \le 0\) the fractional
  ideal \(\mathfrak m_P^{-m} \subset \struct{C,P}\) viewed inside \(K\)).
\end{definition}

\begin{theorem}
[The local step identity]
  \label{mk-thm:adelic_local_step_identity}
  \lean{AlgebraicGeometry.Adelic.sectionOfDivisor_inf_orderGe}
  \leanok
  \dcref{custom}
  \uses{mk-def:adelic_orderGe, mk-def:adelic_sectionOfDivisor}
  Let \(D \le D'\) agree except at a prime divisor \(P \in U\) with
  \(D'(P) = D(P) + 1\). Then
  \[
    \Gamma(U, \struct{C}(D)) \;=\; \Gamma(U, \struct{C}(D')) \cap
      \mathfrak G_P^{\ge -D(P)}.
  \]
\end{theorem}
\begin{proof}
  \leanok
  A nonzero \(f \in \Gamma(U, \struct{C}(D'))\) lies in \(\Gamma(U,
  \struct{C}(D))\) iff its order bound at \(P\) tightens from \(\ge
  -D(P)-1\) to \(\ge -D(P)\) (unchanged at every other point of \(U\),
  where \(D = D'\)), i.e.\ iff \(f \in \mathfrak G_P^{\ge -D(P)}\).
\end{proof}

\begin{theorem}
[The local step quotient injects into the valuation quotient]
  \label{mk-thm:adelic_local_step_injective}
  \lean{AlgebraicGeometry.Adelic.localStepQuot_injective,
    AlgebraicGeometry.Adelic.localStepMapₖ_injective}
  \leanok
  \dcref{custom}
  \uses{mk-thm:adelic_local_step_identity}
  Under the hypotheses of \ref{mk-thm:adelic_local_step_identity}, the
  inclusion \(\Gamma(U, \struct{C}(D')) \subseteq \mathfrak G_P^{\ge
  -D(P)-1}\) induces an injective map on the relative quotients,
  \[
    \Gamma(U, \struct{C}(D')) \big/ \Gamma(U, \struct{C}(D))
      \;\hookrightarrow\;
      \mathfrak G_P^{\ge -D(P)-1} \big/ \mathfrak G_P^{\ge -D(P)}.
  \]
\end{theorem}
\begin{proof}
  \leanok
  Injectivity transports \ref{mk-thm:adelic_local_step_identity} (a
  kernel computation) into the quotient language: a class mapping to
  zero has a representative \(f \in \Gamma(U, \struct{C}(D')) \cap
  \mathfrak G_P^{\ge -D(P)} = \Gamma(U, \struct{C}(D))\), hence is zero
  already in the source quotient.
\end{proof}

\begin{lemma}
[Order-nonnegativity is stalk-integrality]
  \label{mk-lem:adelic_order_stalk_integral}
  \lean{AlgebraicGeometry.Adelic.exists_stalk_lift_of_order_nonneg}
  \leanok
  \dcref{custom}
  \uses{mk-def:adelic_pointValuation}
  For \(f \in K^\times\) with \(v_P(f) \ge 0\), \(f\) lifts to the
  discrete valuation ring \(\struct{C, P}\), i.e.\ \(f = a\) for some
  \(a \in \struct{C, P}\) (viewed in \(K\)).
\end{lemma}
\begin{proof}
  \leanok
  \(v_P(f) \ge 0\) exactly says the normalised valuation of \(f\) at
  \(P\) satisfies \(|f|_P \le 1\); a discrete valuation ring is
  precisely the set of elements of its fraction field of valuation
  \(\le 1\), so \(f\) lifts to \(\struct{C, P}\).
\end{proof}

\begin{theorem}
[The ground field has order zero]
  \label{mk-thm:adelic_constant_field_order_zero}
  \lean{AlgebraicGeometry.Adelic.IsConstantField,
    AlgebraicGeometry.Scheme.isConstantField_functionField}
  \leanok
  \dcref{custom}

  \uses{mk-def:adelic_pointValuation}
  The structure morphism embeds \(k\) in \(K\), and every nonzero
  \(c\in k\) is a unit in every local ring of \(C\).  Hence
  \(v_P(c)=0\) for every prime divisor \(P\).  Consequently every divisor
  section group is naturally a \(k\)-subspace of \(K\).
\end{theorem}

\begin{theorem}
[The local step is bounded by the residue degree]
  \label{mk-thm:adelic_local_step_dim_le_residue_degree}
  \lean{AlgebraicGeometry.Adelic.residueDeg,
    AlgebraicGeometry.Adelic.finrank_localStepTgt,
    AlgebraicGeometry.Adelic.localStep_finrank_le}
  \leanok
  \dcref{custom}
  \uses{mk-thm:adelic_local_step_injective, mk-thm:adelic_constant_field_order_zero,
    mk-def:adelic_pointValuation}
  Put
  \[
    \deg_P^{\mathrm{val}}
      =\dim_k\bigl(\mathfrak G_P^{\ge0}/\mathfrak G_P^{\ge1}\bigr).
  \]
  Assume this degree-zero valuation step is finite-dimensional over \(k\).
  Multiplication by an element of prescribed order identifies every
  quotient \(\mathfrak G_P^{\ge m-1}/\mathfrak G_P^{\ge m}\) with the
  quotient defining \(\deg_P^{\mathrm{val}}\).  Under the hypotheses of
  \ref{mk-thm:adelic_local_step_identity},
  \[
    \dim_k \bigl(\Gamma(U, \struct{C}(D'))/\Gamma(U, \struct{C}(D))\bigr)
      \;\le\; \deg_P^{\mathrm{val}}.
  \]
\end{theorem}
\begin{proof}
  \leanok
  The injection of \ref{mk-thm:adelic_local_step_injective} lands in a
  one-step valuation quotient.  Shifting by a uniformizer power identifies
  this target with the degree-zero step, and dimension is monotone under
  an injective linear map.
\end{proof}

\begin{theorem}[Residue degree one over an algebraically closed base]
  \label{mk-thm:adelic_residue_degree_one}
  \lean{AlgebraicGeometry.Adelic.residueDeg_eq_one_of_isAlgClosed_curve}
  \leanok
  \dcref{custom}
  \uses{mk-def:adelic_residue_degree}
  Let \(C\) be a smooth curve over an algebraically closed field \(k\). Then
  every prime divisor \(P\) of \(C\) satisfies
  \[
    \deg P = [\kappa(P):k] = 1 .
  \]
\end{theorem}
\begin{proof}
  \leanok
  \uses{mk-def:adelic_residue_degree}
  A prime divisor of \(C\) is a closed point, and \(C\) is locally of finite
  type over \(k\) because it is smooth. The underlying space of a scheme
  locally of finite type over a field is a Jacobson space, so
  \(\kappa(P)\) is an algebraic extension of \(k\) by Zariski's lemma; since
  \(k\) is algebraically closed, the structure map \(k \to \kappa(P)\) is an
  isomorphism, giving \(\dim_k \kappa(P) = 1\).
\end{proof}

\begin{corollary}[The weighted degree is the geometric degree]
  \label{mk-cor:adelic_degK_eq_degree}
  \lean{AlgebraicGeometry.Adelic.degK_eq_degree_of_isAlgClosed_curve}
  \leanok
  \dcref{custom}
  \uses{mk-def:adelic_divisor_degree, mk-thm:adelic_residue_degree_one}
  For a smooth curve \(C\) over an algebraically closed field \(k\) and every
  Weil divisor \(D\) on \(C\),
  \[
    \deg_k(D) = \deg(D),
  \]
  the right-hand side being the unweighted degree of \ref{mk-def:divisor_degree}.
\end{corollary}
\begin{proof}
  \leanok
  \uses{mk-thm:adelic_residue_degree_one}
  Both degrees are the sum over the (finite) support of \(D\) of its
  coefficients, weighted by \([\kappa(P):k]\) on the left and by \(1\) on the
  right. The weights agree by \ref{mk-thm:adelic_residue_degree_one}.
\end{proof}

\begin{lemma}[The valuation quotient is the residue field]
  \label{mk-lem:adelic_valuation_quotient_residue_field}
  \dcref{custom}
  \uses{mk-def:adelic_residue_degree, mk-def:adelic_orderGe}
  Reduction modulo the maximal ideal induces a \(k\)-linear isomorphism
  \[
    \mathfrak G_P^{\ge0}/\mathfrak G_P^{\ge1}\simeq_k\kappa(P).
  \]
  In particular \(\deg_P^{\mathrm{val}}=[\kappa(P):k]=\deg P\).
\end{lemma}
\begin{proof}
  The subgroup \(\mathfrak G_P^{\ge0}\) is the valuation ring
  \(\struct{C,P}\), while \(\mathfrak G_P^{\ge1}\) is its maximal ideal.
  Their quotient is therefore the residue field.
\end{proof}

\begin{theorem}[Approximation on a Dedekind chart]
  \label{mk-thm:adelic_weak_approximation}
  \uses{mk-thm:adelic_chartRing_isDedekindDomain,
    mk-def:adelic_pointValuation}
  \dcref{custom}
  Let \(P_1,\ldots,P_r\) be pairwise distinct points in a common affine
  Dedekind chart.  Given \(a_i\in K\) and integers \(m_i\), there exists
  \(f\in K\) such that
  \[
    v_{P_i}(f-a_i)\geq m_i\qquad(1\leq i\leq r).
  \]
\end{theorem}
\begin{proof}
  Write the chart as \(\operatorname{Spec}A\), and choose a common
  denominator \(0\ne d_0\in A\) with \(d_0a_i\in A\) for every \(i\).
  Choose \(N_i\geq0\) with
  \(v_{P_i}(d_0)+N_i+m_i\geq0\).  The nonzero ideal
  \(\prod_i\mathfrak p_i^{N_i}\) contains a nonzero element \(b\); put
  \(d=d_0b\) and \(e_i=v_{P_i}(d)+m_i\).  Then \(da_i\in A\) and
  \(e_i\geq0\).  The powers
  \(\mathfrak p_i^{e_i}\) of the distinct maximal ideals are pairwise
  comaximal.  By the Chinese remainder theorem there is \(y\in A\) with
  \(y\equiv da_i\pmod{\mathfrak p_i^{e_i}}\) for every \(i\).  For
  \(f=y/d\) one has
  \[
    v_{P_i}(f-a_i)=v_{P_i}(y-da_i)-v_{P_i}(d)
      \geq e_i-v_{P_i}(d)=m_i.
  \]
\end{proof}

\begin{theorem}
[The local step has residue degree under approximation]
  \label{mk-thm:adelic_local_step_dim_eq_residue_degree}
  \lean{AlgebraicGeometry.Adelic.localStepMapₖ_surjective,
    AlgebraicGeometry.Adelic.localStep_finrank_eq}
  \leanok
  \dcref{custom}

  \uses{mk-thm:adelic_local_step_injective,
    mk-thm:adelic_local_step_dim_le_residue_degree}
  Suppose every element of \(\mathfrak G_P^{\ge-D(P)-1}\) is congruent,
  modulo \(\mathfrak G_P^{\ge-D(P)}\), to a member of
  \(\Gamma(U,\struct C(D+P))\).  Then
  \[
    \dim_k\bigl(\Gamma(U,\struct C(D+P))/\Gamma(U,\struct C(D))\bigr)
      =\deg_P^{\mathrm{val}}.
  \]
\end{theorem}
\begin{proof}
  \leanok
  The approximation hypothesis makes the local-step map surjective; it is
  injective by \ref{mk-thm:adelic_local_step_injective}.  Hence it is an
  isomorphism onto the one-step valuation quotient.
\end{proof}

\begin{proposition}[Sections of a finite-support quotient]
  \label{mk-thm:adelic_finite_support_quotient_sections}
  \notready
  \dcref{custom}
  \uses{mk-def:adelic_divisorial_sheaf, mk-def:codim1_cycles}
  For \(D\le D'\), the quotient
  \(\mathcal T=\struct C(D')/\struct C(D)\) is supported on the finite set
  \(\operatorname{supp}(D'-D)\), and restriction to stalks induces
  \[
    H^0(C,\mathcal T)\simeq_k
      \bigoplus_{P\in\operatorname{supp}(D'-D)}\mathcal T_P.
  \]
\end{proposition}
\begin{proof}
  Divisors are finitely supported.  The coherent quotient is the pushforward
  of a module on a finite closed subscheme \(Z\subset C\) with this support.
  The finite scheme \(Z\) is affine, and its Artinian coordinate ring is the
  product of its local Artinian factors.  The corresponding module therefore
  decomposes as the direct sum of its localizations at the points of \(Z\),
  which are precisely the stalks \(\mathcal T_P\).  Taking global sections of
  the finite pushforward gives the displayed isomorphism.
\end{proof}

\begin{proposition}[The one-point stalk quotient]
  \label{mk-thm:adelic_twist_stalk_quotient}
  \notready
  \dcref{custom}
  \uses{mk-def:adelic_divisorial_sheaf,
    mk-lem:adelic_valuation_quotient_residue_field}
  If \(D'=D+P\) and
  \(\mathcal T=\struct C(D')/\struct C(D)\), then
  \[
    \mathcal T_P\simeq_k
      \mathfrak G_P^{\ge-D(P)-1}/\mathfrak G_P^{\ge-D(P)}
      \simeq_k\kappa(P),
  \]
  while \(\mathcal T_Q=0\) for \(Q\ne P\).
\end{proposition}
\begin{proof}
  Localising the two divisorial sheaves at \(P\) gives two consecutive
  fractional ideals of the DVR \(\struct{C,P}\).  Multiplication by a
  uniformiser power shifts their quotient to
  \(\struct{C,P}/\mathfrak m_P=\kappa(P)\).  At every other point the two
  divisors, and hence the two stalks, agree.
\end{proof}

\begin{proposition}
[Dimension of the twist quotient]
  \label{mk-thm:adelic_dim_adele_quotient_eq_degree}
  \dcref{custom}
  \uses{mk-thm:adelic_finite_support_quotient_sections,
    mk-thm:adelic_twist_stalk_quotient,
    mk-def:adelic_divisor_degree}
  For \(D \le D'\) the skyscraper section space is finite-dimensional
  with
  \[
    \dim_k \mathcal{Q}(D, D')
      \;=\; \deg_k(D' - D)
      \;=\; \sum_P \bigl(D'(P) - D(P)\bigr)\,\deg P .
  \]
\end{proposition}
\begin{proof}
  By \ref{mk-thm:adelic_finite_support_quotient_sections}, global sections
  split into the finitely many stalk quotients.  Filter each stalk by the
  coefficient difference \(D'(P)-D(P)\).  Every consecutive quotient is
  \(\kappa(P)\) by \ref{mk-thm:adelic_twist_stalk_quotient}, and therefore has
  \(k\)-dimension \(\deg P\).  Additivity of dimension along the filtration
  gives
  \(\dim_k \mathcal{Q}(D, D') = \sum_P (D'(P) - D(P))\deg P = \deg_k(D' - D)\),
  the residue-weighted degree of \ref{mk-def:adelic_divisor_degree}.
\end{proof}

\subsection{Euler characteristic and the Riemann inequality}

\begin{lemma}
[Alternating dimension identity of a four-term exact sequence]
  \label{mk-lem:adelic_finrank_alternating_exact4}
  \lean{AlgebraicGeometry.Adelic.finrank_alternating_of_exact4}
  \leanok
  \dcref{custom}
  For a four-term exact sequence of finite-dimensional \(k\)-vector
  spaces \(0 \to A \xrightarrow{i} B \xrightarrow{p} C \xrightarrow{q} E
  \to 0\) (\(i\) injective, \(\operatorname{im} i = \ker p\),
  \(\operatorname{im} p = \ker q\), \(q\) surjective),
  \[
    \dim A - \dim B + \dim C - \dim E = 0.
  \]
\end{lemma}
\begin{proof}
  \leanok
  Rank--nullity gives \(\dim(\ker p) = \dim(\operatorname{im} i) = \dim
  A\) and \(\dim(\operatorname{im} p) + \dim(\ker p) = \dim B\), so
  \(\dim B - \dim A = \dim(\operatorname{im} p) = \dim(\ker q)\).
  Rank--nullity for \(q\) gives \(\dim(\ker q) + \dim(\operatorname{im}
  q) = \dim C\), and surjectivity gives \(\operatorname{im} q = E\), so
  \(\dim(\ker q) = \dim C - \dim E\). Combining,
  \(\dim B - \dim A = \dim C - \dim E\).
\end{proof}

\begin{definition}
[The \(\chi\)-ledger dimensions]
  \label{mk-def:adelic_chi_ledger}
  \lean{AlgebraicGeometry.Adelic.ell, AlgebraicGeometry.Adelic.h1dim,
    AlgebraicGeometry.Adelic.chi}
  \leanok
  \dcref{custom}
  \uses{mk-def:adelic_H1_window, mk-thm:adelic_constant_field_order_zero}
  Fix a field of constants \(k\). Write \(\ell(D) = \dim_k L(D)\),
  \(h^1(D) = \dim_k \check{H}^1_{\mathrm{ff}}(D)\), and
  \(\chi(D) = \ell(D) - h^1(D)\)
  for the Euler characteristic.
\end{definition}

\begin{theorem}[Global regular functions are constants]
  \label{mk-thm:adelic_global_sections_constants}
  \lean{AlgebraicGeometry.Scheme.globalSectionsAlgEquivBase,
    AlgebraicGeometry.Scheme.unitTopSectionsEquivBase,
    AlgebraicGeometry.Scheme.AffineCoverMVSquare.h0_unit_eq_one}
  \leanok
  \dcref{custom}

  The structure map induces a \(k\)-algebra equivalence
  \[
    \Gamma(C,\struct C)\simeq_k k.
  \]
  Equivalently, the zeroth cohomology of the structure sheaf has dimension
  one on every affine Mayer--Vietoris square.
\end{theorem}
\begin{proof}
  \leanok
  Proper geometric integrality identifies the global section ring with the
  base field; the cover kernel is the global section space.
\end{proof}

\begin{corollary}[The base value of the divisor ledger]
  \label{mk-thm:adelic_ell_zero}
  \dcref{custom}
  \uses{mk-thm:adelic_global_sections_constants,
    mk-def:adelic_divisorial_sheaf, mk-def:adelic_sectionOfDivisor}
  The zero divisorial sheaf is the structure sheaf, its valuation-defined
  global sections satisfy \(L(0)=k\), and therefore \(\ell(0)=1\).
\end{corollary}
\begin{proof}
  The local condition \(v_P(f)\ge0\) at every prime divisor says precisely
  that the rational function \(f\) is regular everywhere on the regular
  curve.  Thus \(L(0)=\Gamma(C,\struct C)\), and
  \ref{mk-thm:adelic_global_sections_constants} identifies this space with
  \(k\).
\end{proof}

\begin{theorem}
[Conditional overlap-ledger identity]
  \label{mk-thm:adelic_ledger_alternating}
  \lean{AlgebraicGeometry.Adelic.ledger_alternating}
  \leanok
  \dcref{custom}
  \uses{mk-lem:adelic_finrank_alternating_exact4,
    mk-thm:adelic_windowMap_injective, mk-def:adelic_chi_ledger}
  Put \(V=U_0\cap U_1\).  Suppose the four displayed spaces are
  finite-dimensional and the window, connecting, and twist maps form an
  exact sequence
  \[
    0\longrightarrow L(D')/L(D)
      \longrightarrow \Gamma(V,\struct C(D'))/\Gamma(V,\struct C(D))
      \longrightarrow\check H^1_{\mathrm{ff}}(D)
      \longrightarrow\check H^1_{\mathrm{ff}}(D')\longrightarrow0.
  \]
  Then
  \[
    \dim_k\bigl(L(D')/L(D)\bigr) - \dim_k\bigl(\Gamma(V,\struct{C}(D'))/
      \Gamma(V,\struct{C}(D))\bigr) + h^1(D) - h^1(D') = 0.
  \]
  For the canonical window map, injectivity follows from
  \ref{mk-thm:adelic_windowMap_injective} when \(D'-D\) is supported on
  \(V\); exactness at the other terms remains a separate hypothesis here.
\end{theorem}
\begin{proof}
  \leanok
  \ref{mk-lem:adelic_finrank_alternating_exact4} applied to the four-term
  sequence, whose injectivity at the left is
  \ref{mk-thm:adelic_windowMap_injective}.
\end{proof}

\begin{theorem}
[One-step \(\chi\)-additivity]
  \label{mk-thm:adelic_chi_additivity}
  \lean{AlgebraicGeometry.Adelic.chi_add,
    AlgebraicGeometry.Adelic.chi_add_eq_residueDeg}
  \leanok
  \dcref{custom}
  \uses{mk-thm:adelic_ledger_alternating,
    mk-thm:adelic_local_step_dim_eq_residue_degree}
  Under the exactness and finite-dimensionality hypotheses of
  \ref{mk-thm:adelic_ledger_alternating} for \(D\le D'\),
  \[
    \chi(D') \;=\; \chi(D) + \dim_k\bigl(\Gamma(V,\struct{C}(D'))/
      \Gamma(V,\struct{C}(D))\bigr).
  \]
  For a one-point bump \(D'=D+P\) with \(P\in V\), the approximation hypothesis of
  \ref{mk-thm:adelic_local_step_dim_eq_residue_degree} gives
  \(\chi(D+P)=\chi(D)+\deg_P^{\mathrm{val}}\).
\end{theorem}
\begin{proof}
  \leanok
  Rank--nullity identifies
  \(\dim_k\bigl(\Gamma(\top,\struct{C}(D'))/\Gamma(\top,\struct{C}(D))\bigr)
  = \ell(D') - \ell(D)\); substituting into
  \ref{mk-thm:adelic_ledger_alternating} and rearranging with
  \(\chi = \ell - h^1\) gives the first claim.  The one-point formula
  follows from \ref{mk-thm:adelic_local_step_dim_eq_residue_degree}.
\end{proof}

\begin{theorem}[Telescoping along an effective divisor]
  \label{mk-thm:adelic_chi_telescope_effective}
  \lean{AlgebraicGeometry.Adelic.pointDivisor,
    AlgebraicGeometry.Adelic.divisorOfList,
    AlgebraicGeometry.Adelic.chi_telescope_list}
  \leanok
  \dcref{custom}

  \uses{mk-thm:adelic_chi_additivity}
  Let \(P_1,\ldots,P_r\) be prime divisors and suppose the one-point
  identity
  \(\chi(E+P)=\chi(E)+\deg_P^{\mathrm{val}}\) holds at every step.  Then
  \[
    \chi(P_1+\cdots+P_r)
      =\chi(0)+\sum_{j=1}^r\deg_{P_j}^{\mathrm{val}}.
  \]
\end{theorem}
\begin{proof}
  \leanok
  Induct on \(r\).  The empty sum is immediate; the induction step applies
  the one-point identity to \(P_1+(P_2+\cdots+P_r)\).
\end{proof}

\begin{theorem}
[Euler characteristic of a divisor]
  \label{mk-thm:adelic_chi_eq_chi_zero_add_degree}
  \uses{mk-thm:adelic_twist_exact_sequence,
    mk-thm:adelic_dim_adele_quotient_eq_degree,
    mk-thm:adelic_module_finite_genus,
    mk-thm:adelic_function_field_H1_zero_compare,
    mk-thm:adelic_ell_zero, mk-def:genus}
  \dcref{custom}
  The Euler characteristic \(\chi(D) = \ell(D) - i(D)\) is additive in the
  degree:
  \[
    \chi(D') - \chi(D) = \deg(D' - D)
    \quad\text{for } D \le D',
    \qquad\text{hence}\qquad
    \chi(D) = \deg D + 1 - g .
  \]
\end{theorem}
\begin{proof}
  By \ref{mk-thm:adelic_module_finite_genus} the base cohomology
  \(\check{H}^1_{\mathrm{ff}}(0)\) is finite-dimensional; the twist sequence
  bound every \(\ell(D), i(D)\) by finite numbers, so all dimensions
  occurring are finite and the alternating-sum identity for an exact
  sequence of finite-dimensional spaces applies. Apply it to the
  four-term sequence of \ref{mk-thm:adelic_twist_exact_sequence} for
  \(D \le D'\):
  \[
    \dim\!\bigl(L(D')/L(D)\bigr) - \dim \mathcal{Q}(D, D')
      + \dim \check{H}^1_{\mathrm{ff}}(D)
      - \dim \check{H}^1_{\mathrm{ff}}(D') = 0 .
  \]
  With \(\dim(L(D')/L(D)) = \ell(D') - \ell(D)\),
  \(\dim \mathcal{Q}(D, D') = \deg(D' - D)\)
  (\ref{mk-thm:adelic_dim_adele_quotient_eq_degree}), and
  \(\dim \check{H}^1_{\mathrm{ff}}(D) = i(D)\), this rearranges to
  \[
    \bigl(\ell(D') - i(D')\bigr) - \bigl(\ell(D) - i(D)\bigr)
      = \deg(D' - D),
    \qquad\text{i.e.}\qquad \chi(D') - \chi(D) = \deg(D' - D).
  \]
  Every divisor \(D\) satisfies \(D - D_- \le D \le D_+\) with
  \(0 \le D_+\); applying the additivity to \(0 \le D_+\) and to
  \(D \le D + D_-\) and adding, one gets \(\chi(D) = \chi(0) + \deg D\)
  for all \(D\) (additivity of \(\chi - \deg\) as a function on the free
  group of divisors, which agrees at \(0\) and is additive on each unit
  step, hence is \(\chi(0)\) plus the homomorphism \(\deg\)). Finally
  \(\chi(0) = \ell(0) - i(0) = 1 - g\): the global sections
  \(L(0)=k\) have dimension \(\ell(0)=1\) by
  \ref{mk-thm:adelic_ell_zero}.  Moreover,
  \(i(0)=\dim_k\check H^1_{\mathrm{ff}}(0)=g\) by
  \ref{mk-thm:adelic_function_field_H1_zero_compare} and
  \ref{mk-thm:adelic_HModuleOne_linearEquiv_cokernel}. Hence
  \(\chi(D) = \deg D + 1 - g\).
\end{proof}

\begin{lemma}
[Nonnegativity of the index of speciality]
  \label{mk-lem:adelic_chi_le_ell}
  \lean{AlgebraicGeometry.Adelic.chi_le_ell}
  \leanok
  \dcref{custom}
  \uses{mk-def:adelic_chi_ledger}
  \(\chi(D) \le \ell(D)\), i.e.\ \(h^1(D) = i(D) \ge 0\).
\end{lemma}
\begin{proof}
  \leanok
  Immediate from \(\chi(D) = \ell(D) - h^1(D)\) and \(h^1(D) \ge 0\).
\end{proof}

\begin{lemma}
[The Riemann inequality from a telescoped \(\chi\)]
  \label{mk-lem:adelic_riemann_inequality_conditional}
  \lean{AlgebraicGeometry.Adelic.riemann_inequality}
  \leanok
  \dcref{custom}
  \uses{mk-lem:adelic_chi_le_ell}
  Given \(\chi(D) = \chi(0) + \deg D\),
  \[
    \deg D + \chi(0) \;\le\; \ell(D).
  \]
\end{lemma}
\begin{proof}
  \leanok
  Substitute \(\chi(D) = \chi(0) + \deg D\) into
  \ref{mk-lem:adelic_chi_le_ell}.
\end{proof}

\begin{corollary}
[Riemann inequality and finiteness of the twist cokernel]
  \label{mk-thm:adelic_riemann_inequality}
  \uses{mk-thm:adelic_chi_eq_chi_zero_add_degree, mk-thm:adelic_module_finite_genus, mk-thm:adelic_twist_exact_sequence, mk-thm:adelic_dim_adele_quotient_eq_degree, mk-lem:adelic_riemann_inequality_conditional}
  \dcref{custom}
  For every Weil divisor \(D\), the space
  \(\check{H}^1_{\mathrm{ff}}(D)\) is
  finite-dimensional and
  \[
    \deg D + 1 - g \;\le\; \ell(D) .
  \]
\end{corollary}
\begin{proof}
  Finiteness: for \(0 \le D''\) with \(D \le D''\) and \(0 \le D''\), the
  last map of the twist sequence
  (\ref{mk-thm:adelic_twist_exact_sequence}) makes
  \(\check{H}^1_{\mathrm{ff}}(D) \to
  \check{H}^1_{\mathrm{ff}}(D'')\) surjective with kernel a
  quotient of \(\mathcal{Q}(D, D'')\), which is finite
  (\ref{mk-thm:adelic_dim_adele_quotient_eq_degree}); and
  \(\check{H}^1_{\mathrm{ff}}(0) \to
  \check{H}^1_{\mathrm{ff}}(D'')\) is surjective, so
  \(\dim \check{H}^1_{\mathrm{ff}}(D'')
    \le \dim \check{H}^1_{\mathrm{ff}}(0) = g\). Hence
  \(\dim \check{H}^1_{\mathrm{ff}}(D)
    \le g + \deg(D'' - D) < \infty\). Inequality:
  \ref{mk-lem:adelic_riemann_inequality_conditional} specialises, with
  \(\chi(0) = 1-g\)
  (\ref{mk-thm:adelic_chi_eq_chi_zero_add_degree}), to
  \(\ell(D) = \deg D + 1 - g + i(D) \ge \deg D + 1 - g\).
\end{proof}

\section{Weil differentials and duality}
\label{mk-sec:RiemannRoch_Adelic_duality}

The adelic evaluation pairing identifies the dual of the twist cokernel
with a space of Weil differentials.  A fixed nonzero differential then
turns this space into an ordinary Riemann--Roch space.

\begin{proposition}[The principal-parts resolution]
  \label{mk-thm:adelic_principal_parts_resolution}
  \notready
  \dcref{custom}
  \uses{mk-def:adelic_divisorial_sheaf, mk-def:adelic_orderGe}
  Let \(\mathcal K\) be the sheaf of rational functions and let
  \[
    \mathcal P_D=
      \bigoplus_P i_{P,*}
        \bigl(K/\mathfrak G_P^{\ge-D(P)}\bigr).
  \]
  There is an exact sequence
  \[
    0\longrightarrow\struct C(D)\longrightarrow\mathcal K
      \longrightarrow\mathcal P_D\longrightarrow0,
  \]
  in which \(\mathcal K\) and \(\mathcal P_D\) are flasque, and
  \(\Gamma(C,\mathcal P_D)\simeq
    \mathbb A_K/\mathbb A_K(D)\).
\end{proposition}
\begin{proof}
  Exactness is stalkwise.  At the generic point the first arrow is the
  identity of \(K\).  At a closed point \(P\), its kernel is the
  fractional ideal \(\mathfrak G_P^{\ge-D(P)}\), and the quotient is the
  displayed principal-parts module.  The rational-function sheaf is
  constant on nonempty opens of the integral curve and is flasque; each
  skyscraper summand is flasque, and a finite-support section of their
  direct sum extends componentwise.  Finally an adele modulo
  \(\mathbb A_K(D)\) records exactly its finitely many nonzero principal
  parts, which gives the global-section isomorphism.
\end{proof}

\begin{proposition}[Adelic and two-chart cohomology agree]
  \label{mk-thm:adelic_adele_cover_compare}
  \notready
  \dcref{custom}
  \uses{mk-thm:adelic_principal_parts_resolution,
    mk-thm:adelic_function_field_H1_compare}
  For every divisor \(D\),
  \[
    H^1_{\mathbb A}(D)
      =\mathbb A_K/(\mathbb A_K(D)+K)
      \simeq_k\check H^1_{\mathrm{ff}}(D).
  \]
\end{proposition}
\begin{proof}
  Taking global sections of
  \ref{mk-thm:adelic_principal_parts_resolution} identifies its degree-one
  cohomology with the cokernel of
  \(K\to\mathbb A_K/\mathbb A_K(D)\), namely
  \(\mathbb A_K/(\mathbb A_K(D)+K)\).  Thus this quotient is
  \(H^1(C,\struct C(D))\), which is the two-chart quotient by
  \ref{mk-thm:adelic_function_field_H1_compare}.
\end{proof}

\begin{definition}[Weil differentials]
  \label{mk-def:adelic_weil_differentials}
  \uses{mk-thm:adelic_adele_cover_compare}
  \dcref{custom}
  A Weil differential is a \(k\)-linear form on \(\mathbb A_K\) that
  vanishes on \(\mathbb A_K(E)+K\) for some divisor \(E\).  Write
  \(\Omega_K\) for their union and
  \[
    \Omega_K(D)=
      \{\eta\in\Omega_K:\eta(\mathbb A_K(D)+K)=0\}.
  \]
  Multiplication by \(x\in K\) is defined by
  \((x\eta)(\alpha)=\eta(x\alpha)\).
\end{definition}
\begin{proof}[Closure under the vector-space operations]
  If \(\eta\) vanishes on \(\mathbb A_K(E)+K\) and \(\theta\) vanishes
  on \(\mathbb A_K(F)+K\), then both vanish on
  \(\mathbb A_K(\min(E,F))+K\), where the minimum is taken
  coefficientwise.  Hence so does \(\eta+\theta\).  If \(x\ne0\), then
  \[
    x\mathbb A_K(E+\operatorname{div}(x))=\mathbb A_K(E)
    \quad\text{and}\quad xK=K,
  \]
  so \((x\eta)(\alpha)=\eta(x\alpha)\) vanishes on
  \(\mathbb A_K(E+\operatorname{div}(x))+K\).  The zero scalar is
  immediate.
  Thus the displayed operation makes \(\Omega_K\) a \(K\)-vector
  space.
\end{proof}

\begin{proposition}[The Weil-differential pairing]
  \label{mk-thm:adelic_residuePairing}
  \notready
  \source{papaioannou-algebraic-rr:page-0011}
  \uses{mk-def:adelic_weil_differentials,
    mk-thm:adelic_adele_cover_compare}
  \dcref{custom}
  Evaluation induces a canonical isomorphism
  \[
    \Omega_K(D)\simeq_k H^1_{\mathbb A}(D)^*,
  \]
  and hence a perfect pairing
  \[
    \check H^1_{\mathrm{ff}}(D)\times\Omega_K(D)\longrightarrow k.
  \]
\end{proposition}
\begin{proof}
  A functional in \(\Omega_K(D)\) factors uniquely through the
  finite-dimensional quotient
  \(H^1_{\mathbb A}(D)=\mathbb A_K/(\mathbb A_K(D)+K)\), and every linear
  functional on that quotient pulls back to such a Weil differential.
  This gives the first isomorphism; transport it across
  \ref{mk-thm:adelic_adele_cover_compare} for the displayed pairing.
\end{proof}

\begin{theorem}[Eventual vanishing of the adelic index]
  \label{mk-thm:adelic_index_eventually_zero}
  \notready
  \dcref{custom}
  \uses{mk-thm:adelic_module_finite_genus,
    mk-thm:adelic_adele_cover_compare,
    mk-thm:adelic_riemann_inequality,
    mk-thm:adelic_principal_degree_zero}
  There is an integer \(c\) such that
  \(H^1_{\mathbb A}(D)=0\) whenever \(\deg D\ge c\).
\end{theorem}
\begin{proof}
  Choose adele representatives of a basis of the finite-dimensional space
  \(H^1_{\mathbb A}(0)\).  Every representative belongs to
  \(\mathbb A_K(B_j)\) for some effective divisor \(B_j\); a common upper
  bound \(B\) therefore kills the basis under the surjection
  \(H^1_{\mathbb A}(0)\to H^1_{\mathbb A}(B)\), so the latter space is zero.
  Put \(c=\deg B+g\).  If \(\deg D\ge c\), the Riemann inequality gives
  \(\ell(D-B)>0\).  For nonzero \(f\in L(D-B)\), the divisor
  \(E=D+\operatorname{div}(f)\) satisfies \(E\ge B\), hence
  \(H^1_{\mathbb A}(E)=0\).  Multiplication by \(f^{-1}\) identifies the
  quotients for \(D\) and \(E\), proving the claim.
\end{proof}

\begin{theorem}[The divisor of a Weil differential]
  \label{mk-thm:adelic_weil_differential_divisor}
  \notready
  \source{papaioannou-algebraic-rr:page-0012}
  \uses{mk-def:adelic_weil_differentials,
    mk-thm:adelic_residuePairing,
    mk-thm:adelic_index_eventually_zero}
  \dcref{custom}
  For every nonzero \(\eta\in\Omega_K\) there is a unique divisor
  \(\operatorname{div}(\eta)\) maximal among the divisors \(D\) for which
  \(\eta\in\Omega_K(D)\).  Consequently
  \[
    \Omega_K(D)=\{0\}\cup
      \{\eta\ne0:\operatorname{div}(\eta)\ge D\}.
  \]
\end{theorem}
\begin{proof}
  The set \(M(\eta)=\{D:\eta\in\Omega_K(D)\}\) is nonempty by the
  definition of a Weil differential.  If \(D\in M(\eta)\), then \(\eta\)
  gives a nonzero functional on \(H^1_{\mathbb A}(D)\); hence
  \ref{mk-thm:adelic_index_eventually_zero} bounds the degrees occurring in
  \(M(\eta)\).  Choose \(W\in M(\eta)\) of maximal degree.  If
  \(D_0\in M(\eta)\) and \(D_0(Q)>W(Q)\), decompose any
  \(\alpha\in\mathbb A_K(W+Q)\) into the adele supported away from \(Q\)
  and the adele supported at \(Q\).  The first lies in \(\mathbb A_K(W)\)
  and the second in \(\mathbb A_K(D_0)\), so \(\eta(\alpha)=0\).  Thus
  \(W+Q\in M(\eta)\), contradicting maximality.  Therefore every
  \(D\in M(\eta)\) satisfies \(D\le W\), which proves existence,
  uniqueness, and the displayed characterization.
\end{proof}

\begin{theorem}[Scaling a differential]
  \label{mk-thm:adelic_weil_differential_divisor_mul}
  \notready
  \source{papaioannou-algebraic-rr:page-0012}
  \uses{mk-thm:adelic_weil_differential_divisor}
  \dcref{custom}
  For \(x\in K^\times\) and \(0\ne\eta\in\Omega_K\),
  \[
    \operatorname{div}(x\eta)
      =\operatorname{div}(x)+\operatorname{div}(\eta).
  \]
\end{theorem}
\begin{proof}
  Multiplication by \(x\) carries \(\mathbb A_K(E)\) isomorphically onto
  \(\mathbb A_K(E-\operatorname{div}(x))\).  Hence
  \(x\eta\in\Omega_K(E)\) exactly when
  \(\eta\in\Omega_K(E-\operatorname{div}(x))\).  Comparing the unique
  maximal divisors on the two sides gives the formula.
\end{proof}

\begin{theorem}[One-dimensionality of Weil differentials]
  \label{mk-thm:adelic_weil_differentials_finrank_one}
  \notready
  \source{papaioannou-algebraic-rr:page-0011,
    papaioannou-algebraic-rr:page-0012}
  \uses{mk-thm:adelic_residuePairing,
    mk-thm:adelic_weil_differential_divisor,
    mk-thm:adelic_weil_differential_divisor_mul,
    mk-thm:adelic_riemann_inequality,
    mk-thm:adelic_negative_degree_vanishing,
    mk-thm:adelic_chi_eq_chi_zero_add_degree}
  \dcref{custom}
  The \(K\)-vector space \(\Omega_K\) is one-dimensional.
\end{theorem}
\begin{proof}
  First \(\Omega_K\ne0\): for \(\deg D\ll0\), the identity
  \(i(D)=\ell(D)-\deg D-1+g\) makes \(i(D)>0\), so
  \(\Omega_K(D)\simeq H^1_{\mathbb A}(D)^*\) is nonzero.  Choose
  \(0\ne\eta_i\in\Omega_K(D_i)\) for \(i=1,2\).  For an effective
  divisor \(B\) of sufficiently large degree, multiplication gives
  injections
  \[
    L(D_i+B)\hookrightarrow\Omega_K(-B),\qquad x\longmapsto x\eta_i.
  \]
  The Riemann inequality bounds the two source dimensions below by
  \(\deg(D_i+B)+1-g\), while negative-degree vanishing and the
  \(\chi\)-identity give
  \(\dim_k\Omega_K(-B)=i(-B)=\deg B+g-1\).  For \(\deg B\) large, the
  two images therefore have positive-dimensional intersection.  Thus
  \(x_1\eta_1=x_2\eta_2\ne0\) for some \(x_i\in K^\times\), and
  \(\eta_2=(x_1/x_2)\eta_1\).  Every two nonzero differentials are
  proportional.
\end{proof}

\begin{theorem}
[Serre duality: \(i(D) = \ell(K_C - D)\)]
  \label{mk-thm:adelic_serreDuality}
  \notready
  \source{papaioannou-algebraic-rr:page-0011,
    papaioannou-algebraic-rr:page-0012,
    papaioannou-algebraic-rr:page-0013}
  \uses{mk-thm:adelic_residuePairing,
    mk-thm:adelic_adele_cover_compare,
    mk-thm:adelic_weil_differentials_finrank_one,
    mk-thm:adelic_weil_differential_divisor_mul}
  \dcref{custom}
  For a canonical divisor \(K_C = \operatorname{div}\omega\) of a nonzero
  Weil differential \(\omega\), the map \(x \mapsto x\omega\) induces a
  \(k\)-linear isomorphism
  \[
    \mu \colon L(K_C - D) \;\xrightarrow{\ \sim\ }\; \Omega_K(D)
      \;\simeq_k\; H^1_{\mathbb A}(D)^*
      \;\simeq_k\;\check H^1_{\mathrm{ff}}(D)^*,
  \]
  so that \(i(D)=\dim_k\check H^1_{\mathrm{ff}}(D)=\ell(K_C-D)\).
\end{theorem}
\begin{proof}
  By \ref{mk-thm:adelic_weil_differentials_finrank_one}, every Weil
  differential is uniquely \(x\omega\).  By
  \ref{mk-thm:adelic_weil_differential_divisor_mul},
  \(\operatorname{div}(x\omega)=\operatorname{div}(x)+K_C\); hence
  \[
    x\omega\in\Omega_K(D)
      \quad\Longleftrightarrow\quad
    \operatorname{div}(x)+K_C\ge D
      \quad\Longleftrightarrow\quad x\in L(K_C-D).
  \]
  This proves the first isomorphism.  The second is
  \ref{mk-thm:adelic_residuePairing}, and the third is
  \ref{mk-thm:adelic_adele_cover_compare}; taking dimensions gives the final
  equality.
\end{proof}

\begin{theorem}[Geometric degree of a principal divisor, from the \(\chi\)-ledger]
  \label{mk-thm:adelic_principal_geometric_degree_zero}
  \lean{AlgebraicGeometry.Adelic.degree_principal_eq_zero_of_isAlgClosed_curve}
  \leanok
  \dcref{custom}
  \uses{mk-cor:adelic_degK_eq_degree, mk-thm:adelic_chi_eq_chi_zero_add_degree}
  Let \(C\) be a smooth curve over an algebraically closed field \(k\), and
  suppose the \(\chi\)-ledger
  \(\chi(D)=\chi(0)+\deg_k(D)\) holds at every Weil divisor \(D\). Then for
  every \(g\in K^\times\),
  \[
    \deg\bigl(\operatorname{div}(g)\bigr)=0 .
  \]
\end{theorem}
\begin{proof}
  \leanok
  \uses{mk-cor:adelic_degK_eq_degree}
  The invariant \(\chi\) is unchanged along a linear equivalence, so
  \(\chi(\operatorname{div}(g))=\chi(0)\); the ledger then forces
  \(\deg_k(\operatorname{div}(g))=0\). Over an algebraically closed base the
  weighted and geometric degrees coincide by
  \ref{mk-cor:adelic_degK_eq_degree}, giving the displayed equality.
\end{proof}

\begin{lemma}[Negative-degree divisors have no nonzero sections]
  \label{mk-thm:adelic_negative_degree_vanishing}
  \uses{mk-thm:adelic_principal_degree_zero,
    mk-def:adelic_sectionOfDivisor}
  \dcref{custom}
  If \(\deg E<0\), then \(L(E)=0\) and \(\ell(E)=0\).
\end{lemma}
\begin{proof}
  A nonzero \(f\in L(E)\) would make
  \(\operatorname{div}(f)+E\) effective.  Its degree is \(\deg E<0\),
  because principal divisors have weighted degree zero by
  \ref{mk-thm:adelic_principal_degree_zero}; an effective divisor cannot
  have negative degree.
\end{proof}

\begin{theorem}[Degree of a canonical divisor]
  \label{mk-thm:adelic_canonical_degree}
  \notready
  \source{papaioannou-algebraic-rr:page-0013}
  \uses{mk-thm:adelic_serreDuality,
    mk-thm:adelic_chi_eq_chi_zero_add_degree,
    mk-thm:adelic_ell_zero}
  \dcref{custom}
  Every canonical divisor \(K_C\) has
  \[
    \ell(K_C)=g,
    \qquad \deg K_C=2g-2.
  \]
\end{theorem}
\begin{proof}
  At \(D=0\), duality gives \(i(0)=\ell(K_C)\), while
  \(\ell(0)=1\), \(i(0)=g\), so \(\ell(K_C)=g\).  At \(D=K_C\),
  duality gives \(i(K_C)=\ell(0)=1\).  Substituting these two dimensions
  into \(\ell(K_C)-i(K_C)=\deg K_C+1-g\) yields
  \(\deg K_C=2g-2\).
\end{proof}

\begin{corollary}
[Vanishing: \(i(D) = 0\) for \(\deg D > 2g - 2\)]
  \label{mk-thm:adelic_h1_vanishing}
  \notready
  \source{papaioannou-algebraic-rr:page-0013}
  \uses{mk-thm:adelic_serreDuality,
    mk-thm:adelic_canonical_degree,
    mk-thm:adelic_negative_degree_vanishing,
    mk-thm:adelic_function_field_H1_compare}
  \dcref{custom}
  If \(\deg D>2g-2\), then
  \(\check H^1_{\mathrm{ff}}(D)=0\), and consequently
  \(H^1(C,\struct C(D))=0\).  This is the curve-level replacement for
  Castelnuovo--Mumford boundedness in the Quot endgame.
\end{corollary}
\begin{proof}
  By \ref{mk-thm:adelic_canonical_degree}, \(\deg(K_C-D)<0\).  Therefore
  \ref{mk-thm:adelic_negative_degree_vanishing} and duality give
  \(i(D)=\ell(K_C-D)=0\).  Finally
  \ref{mk-thm:adelic_function_field_H1_compare} supplies the sheaf-cohomology
  statement.
\end{proof}

\begin{theorem}
[Riemann--Roch]
  \label{mk-thm:adelic_riemannRoch}
  \notready
  \source{papaioannou-algebraic-rr:page-0013}
  \uses{mk-thm:adelic_chi_eq_chi_zero_add_degree, mk-thm:adelic_serreDuality}
  \dcref{custom}
  For every Weil divisor \(D\) and canonical divisor \(K_C\),
  \[
    \ell(D) - \ell(K_C - D) \;=\; \deg D + 1 - g .
  \]
\end{theorem}
\begin{proof}
  Immediate from the two keystones: the \(\chi\)-ledger
  \ref{mk-thm:adelic_chi_eq_chi_zero_add_degree} gives
  \(\ell(D) - i(D) = \deg D + 1 - g\), and Serre duality
  \ref{mk-thm:adelic_serreDuality} gives \(i(D) = \ell(K_C - D)\);
  substituting the second into the first yields the identity.
\end{proof}
